\chapter{Compactness and decreasing regularization}
\index{compactness!BFJ psh functions}
\index{regularization!decreasing BFJ}
\index{equicontinuity!on dual complexes}
\label{chap:bfj-compactness}

Two theorems turn the skeletal definition into an analytic theory.
Normalized psh functions form a compact space, and every singular psh
function can be approached monotonically from above by model psh functions.
Their proofs have different engines.  Compactness comes from relative
intersection bounds and toroidal subdivisions; regularization comes from
envelopes and multiplier ideals.

The difference between the two engines is worth stating before either is
run, because it governs which hypotheses are needed where.  The compactness
theorem is, at bottom, an inequality between intersection numbers on one
fixed model, propagated along the incidence graph of a special fibre.  It
uses no vanishing theorem, no resolution beyond the existence of the SNC
model one starts from, and no ideal sheaf.  Once the inequality is in hand,
the topology does the rest: Arzel\`a--Ascoli on each dual complex, Tychonoff
on the product over all of them.  The regularization theorem is of a
different nature.  It asserts that an arbitrary upper semicontinuous
object, possibly with a polar set, is the limit of finite piecewise affine
data.  No inequality between fixed intersection numbers can produce such a
statement; what produces it is the continuity of psh envelopes, and the only
known proof of that continuity in the present generality runs through
asymptotic multiplier ideals on models.  That is where residue
characteristic zero enters, and \Cref{sec:ch22-char-zero} says exactly
where and exactly how far the statement is presently known to extend.

A second theme runs through the chapter and is, if anything, more likely to
be mishandled than the algebra.  The index set of the theory is the system
of models, and that system is directed but in general has no countable
cofinal subsystem.  Compactness of the normalized psh space is therefore a
statement about \emph{nets}: every net has a convergent subnet, and the
regularizing family is a decreasing net, not a decreasing sequence.  We keep
this visible at every step, use the net toolkit proved in
\Cref{chap:topological-tools}, and mark each place where a sequential
formulation would require a hypothesis that has not been proved.  The
distinction is not pedantic: \Cref{ex:ch22-countable-blind-spot} exhibits two
psh functions at uniform distance one from each other whose restrictions to
every skeleton in a prescribed countable family coincide, so that a diagonal
extraction over countably many skeleta cannot even identify a limit.

\paragraph{Standing hypotheses.}
We retain the original BFJ hypotheses of \Cref{chap:bfj-model-forms}: the
field \(K\) is complete and discretely valued with residue characteristic
zero and \(v(\pi)=1\); the variety \(X\) is smooth, connected and projective
of dimension \(n\) over \(K\); models are normal, flat and projective over
\(R=K^\circ\).  We retain the form \(\theta\) of
\eqref{eq:ch21-standing-theta}, with ample generic class \(\{\theta\}\), and
the class \(\PSH(X,\theta)\) of \Cref{def:ch21-bfj-psh}.  On each face of
every dual complex, choose once and for all a Euclidean norm compatible with
its rational affine structure in the sense of
\Cref{sec:ch22-moduli}.  Different choices change the constants below but
not the statements.

\paragraph{What the reader should learn.}
The reader should understand how intersection inequalities bound vertex
values and slopes, with constants that can be computed on an example; why
compactness uses products and nets rather than a diagonal sequence; what
the multiplier-ideal package on models actually asserts and where it is
used; and exactly where characteristic zero enters decreasing
regularization.  The prerequisites are
\Cref{chap:bfj-model-forms,chap:bfj-psh}, the skeletal material of
\Cref{chap:skeleta} --- dual complexes, retractions, integral affine
structures, projective subdivisions, and cofinality of SNC models --- and
the material on nets, Dini's theorem, upper semicontinuous regularization
and Choquet selection in \Cref{chap:topological-tools}, summarized in
\Cref{app:nets}.  The multiplier and test ideal algebra is collected in
\Cref{app:ideals}; \Cref{sec:ch22-multiplier} states precisely which of its
statements are used here and in which form.

\section{Orientation: two theorems, two engines}
\label{sec:ch22-orientation}

\subsection{The two statements in advance}
\label{sec:ch22-statements-preview}

The chapter proves the two implications

\[
 \begin{aligned}
 &\text{intersection inequalities on one model}\\[-1mm]
 &\qquad\Longrightarrow
   \text{equicontinuity on every dual complex}\\[-1mm]
 &\qquad\Longrightarrow
   \text{compactness of }\PSH(X,\theta)/\mathbf R,\\[1mm]
 &\text{continuity of psh envelopes}\\[-1mm]
 &\qquad\Longrightarrow
   \text{pointwise decreasing model regularization},
 \end{aligned}
\tag{22.1}\label{eq:ch22-architecture}
\]

and organizes their inputs so that no later chapter can attribute to one
engine what belongs to the other.  In detail, the outcome is the following.

\begin{enumerate}
\item \emph{Skeletal estimate.}  For every SNC model \(\mathcal X\) on which
\(\theta\) is determined there is a constant \(C(\mathcal X,\theta)\) such
that the restriction to \(\Delta_{\mathcal X}\) of every \(\theta\)-psh
function is \(C\)-Lipschitz on each face, and lies in \([-C,0]\) once
normalized by \(\sup\varphi=0\).  This is
\Cref{thm:bfj-skeletal-equicontinuity}, and the constant is made explicit in
\Cref{prop:ch22-vertex-estimate,prop:ch22-slope-estimate}.
\item \emph{Compactness.}  The normalization map is proper for the topology
of uniform convergence on all dual complexes, and the normalized space
\(\PSH_0(X,\theta)\) is compact; equivalently
\(\PSH(X,\theta)/\mathbf R\) is compact.  This is
\Cref{thm:ch22-proper-normalization} and
\Cref{thm:bfj-psh-compactness}, in net form throughout.
\item \emph{Decreasing regularization.}  Every \(\varphi\in\PSH(X,\theta)\)
is the pointwise limit of a decreasing net of \(\theta\)-psh model
functions, and a continuous \(\varphi\) is a uniform limit of a decreasing
sequence of them.  This is
\Cref{thm:bfj-decreasing-regularization} and
\Cref{cor:ch22-continuous-decreasing-sequence}.
\end{enumerate}

\subsection{Why the two proofs cannot be merged}
\label{sec:ch22-two-engines}

It is tempting to hope that compactness alone yields regularization: the
model psh functions dominating \(\varphi\) form a net in a compact space, so
some subnet converges, and one only has to identify the limit.  The hope
fails at the identification step, and it fails for a reason that is
structural rather than technical.  Compactness produces a limit in
\(\PSH(X,\theta)\); it says nothing about \emph{which} psh function the
limit is.  To force the limit to be \(\varphi\) one must know that model psh
majorants come arbitrarily close to \(\varphi\) at each prescribed point,
which is exactly the assertion that the psh envelope of a continuous
obstacle is approximable by model functions.  That assertion is the
multiplier-ideal theorem \Cref{thm:bfj-envelope-regularity}, proved in
\Cref{chap:bfj-envelopes}.

The converse temptation is equally common: to hope that regularization makes
compactness elementary, since a decreasing net of model functions is
determined by finitely many numbers at each stage.  This also fails: the
numbers in question are the vertex values and slopes of the approximants,
and it is precisely the uniform bound on those numbers, valid for all
approximants simultaneously, that the skeletal estimate provides.  Without
it, a decreasing net of model functions may have infimum \(-\infty\), and
\Cref{thm:ch21-decreasing-stability} would be vacuous.

The two engines are therefore genuinely independent, and we run them in the
order dictated by their dependencies: the intersection estimate first,
because \Cref{chap:bfj-envelopes} uses it; the order-theoretic
consequences of envelope regularity afterwards, with the envelope theorem
quoted with its exact hypotheses.

\subsection{Hypothesis signature of the chapter}
\label{sec:ch22-signature}

The margin code of \Cref{chap:bfj-model-forms} applies to every statement
below.  The following table records what each block of the chapter actually
consumes; the last column is the one that changes when the ground field is
changed, and it is the subject of \Cref{sec:ch22-char-zero}.

\begin{center}
\small
\begin{tabular}{p{0.24\textwidth}p{0.30\textwidth}p{0.36\textwidth}}
\toprule
Block & Geometric input & Sensitive hypothesis\\
\midrule
Vertex estimate
 & relative nefness, connectedness of \(\mathcal X_s\),
   projection formula
 & none beyond existence of one SNC model\\
Slope estimate
 & projective rational subdivisions, vertical
   \(\mathbf Q\)-factoriality, negativity lemma
 & none beyond existence of one SNC model\\
Equicontinuity for all psh functions
 & the two estimates and uniform approximation on a fixed complex
 & none\\
Compactness
 & Arzel\`a--Ascoli, Tychonoff, the estimates
 & none; but the index set is uncountable\\
Definition of the topology
 & the inverse system of dual complexes
 & cofinality of SNC models, hence resolution\\
Envelope regularity
 & log resolutions, asymptotic multiplier ideals, subadditivity,
   uniform global generation
 & residue characteristic zero\\
Decreasing net
 & order theory, Dini
 & none, once envelope regularity is available\\
\bottomrule
\end{tabular}
\end{center}

Two rows deserve immediate comment.  First, the two intersection estimates
are characteristic free \emph{as estimates}: they are proved on a given SNC
model, using a toroidal blowup that exists in any residue characteristic by
\Cref{thm:skeleton-subdivision}.  What is not characteristic free is the
supply of SNC models, that is, their cofinality
(\Cref{prop:snc-cofinal,rem:cofinality-scope}), and it is cofinality that
makes the collection of dual complexes large enough to define the topology
and to reconstruct a psh function from its skeletal restrictions.  This is
isolated as \Cref{prop:ch22-characteristic-free-estimates}.  Second, the
uncountability in the compactness row is not an artifact of the proof: it is
the content of \Cref{ex:ch22-no-countable-cofinal}.

\subsection{What is imported}
\label{sec:ch22-imports}

Exactly four external statements are used in this chapter and not proved in
it.  Each is quoted where it is used, and each is attributed.

\begin{enumerate}
\item \emph{Realization of projective subdivisions.}  A rational simplicial
projective subdivision of \(\Delta_{\mathcal X}\) is realized by a vertical
blowup with normal, vertically \(\mathbf Q\)-factorial total space; this is
\Cref{thm:skeleton-subdivision}, proved in \Cref{chap:skeleta} after
\cite[Theorem~3.15]{BFJ16}.
\item \emph{Cofinality of SNC models}, \Cref{prop:snc-cofinal}, which rests
on resolution of singularities in the form recorded as
\Cref{thm:appB-resolution}.
\item \emph{Envelope regularity}, \Cref{thm:bfj-envelope-regularity}, proved
in \Cref{chap:bfj-envelopes} from the multiplier-ideal package on models of
\cite[Appendix~B]{BFJ16}; \Cref{sec:ch22-multiplier} states that package in
the exact form consumed here.
\item \emph{Kleiman's criterion} and the projection formula for relative
intersection numbers, in the form of
\Cref{thm:appF-projection,cor:appF-generically-finite}.
\end{enumerate}

Everything else, including the negativity lemma, the simplex lemma, the
Arzel\`a--Ascoli and Tychonoff arguments in the form used, the
order-theoretic construction of the regularizing net, and the continuous
majorant lemma, is proved here.

\section{Normalization, metrics, and moduli on dual complexes}
\label{sec:ch22-moduli}

\subsection{The normalization}

Adding a constant changes no slope, and the class \(\PSH(X,\theta)\) is
stable under addition of constants by
\Cref{prop:ch21-elementary-stability}.  We therefore normalize by

\[
                         \sup_{X^{\An}}\varphi=0.
\tag{22.2}\label{eq:ch22-normalization}
\]

By \eqref{eq:ch21-sup-on-skeleton}, this maximum is visible on every dual
complex on which \(\theta\) is determined: it is attained, and attained
already on \(\Delta_{\mathcal X}\).  Two consequences are used constantly.
First, \(\varphi\leq0\) everywhere, so all the quantities estimated below
have a sign.  Second, the normalization is \emph{skeletal}: it can be tested
on any one determining complex, and the test is stable under enlarging the
complex.  A normalization by the value at a single divisorial point would
not have the second property, and would produce a family that is bounded on
one skeleton and unbounded on another.

\subsection{Integral affine structure, lattices, and norms}

Let \(\mathcal X\) be an SNC model and \(\sigma\) a face of
\(\Delta_{\mathcal X}\), say of dimension \(q\).  By
\Cref{def:integral-affine,prop:integral-affine-intrinsic}, the face carries
an intrinsic integral affine structure: the lattice
\(\operatorname{Aff}_{\mathbf Z}(\sigma)\) of integral affine functions on
\(\sigma\) is well defined, independently of the model realizing
\(\sigma\), and by \Cref{prop:affine-structure-refinement} it is unchanged
under refinement.  Its lattice of affine-linear parts is a lattice
\(M_\sigma\simeq\mathbf Z^q\) in the dual \(V_\sigma^\vee\) of the direction
space \(V_\sigma\) of \(\sigma\), and we write
\(N_\sigma\subset V_\sigma\) for the dual lattice.

\begin{definition}[Compatible norm]
\label{def:ch22-compatible-norm}
A Euclidean norm \(|\cdot|_\sigma\) on \(V_\sigma\) is \emph{compatible with
the integral affine structure} if some basis of \(N_\sigma\) is orthonormal
for it.  A \emph{system of compatible norms} on \(\Delta_{\mathcal X}\) is a
choice of such a norm on every face.
\end{definition}

Compatible norms exist on each face by definition of a lattice.  We do
\emph{not} require the norm of a face to restrict to the norm of a subface,
and in general no system has that property: on the standard triangle with
lattice \(\mathbf Z^2\), the primitive vector \((1,-1)\) of the lattice of
the diagonal edge has Euclidean length \(\sqrt2\), so the compatible norm of
the edge and the restricted norm of the triangle differ.  Nothing below
needs the two to agree: the facewise Lipschitz constant
\eqref{eq:ch22-facewise-lipschitz} is computed face by face, and any two
systems of norms on the finitely many faces of a fixed complex differ by
bounded factors.  This is the meaning of the sentence in the standing
hypotheses: different choices change the constants, not the statements.

What is stable, and is used, is the behaviour under refinement.  By
\Cref{prop:affine-structure-refinement}, subdividing a face does not change
its integral affine structure, so a compatible norm on a face restricts to a
compatible norm on each cell of a subdivision of it.  In dimension one this
is the lattice length of \Cref{rem:lattice-length}: an edge whose endpoints
are the divisorial points of components of multiplicities \(b,b'\) has
lattice length \(1/\operatorname{lcm}(b,b')\), and the compatible norm is
the one for which that edge has length equal to its lattice length.

\begin{example}[Two normalizations of one edge]
\label{ex:ch22-edge-normalization}
For the node chart \(xy=\pi\) of \Cref{ex:refinement-segment}, the edge is
parametrized by \(t=v(x)\in[0,1]\) and the integral affine functions are
\(ct+d\) with \(c,d\in\mathbf Z\); the compatible norm is \(|\rd t|\), and the
edge has length \(1\).  After the single blowup of the node, the two
half-edges have lattice length \(\tfrac12\) each and the parameter \(t\) is
still the compatible coordinate, so the compatible norms of the refined
complex glue to the original one.  Had one instead declared each half-edge
to have length \(1\), the Lipschitz constants of one and the same function
would drop by a factor of two upon refinement, and no uniform estimate could
be stated.
\end{example}

\subsection{Facewise Lipschitz functions and the path metric}

\begin{definition}[Facewise Lipschitz constant]
\label{def:ch22-facewise-lipschitz}
For \(u:\Delta_{\mathcal X}\to\mathbf R\) put

\[
 \operatorname{Lip}_{\mathrm f}(u)
 :=\max_{\sigma}\ \sup_{x\neq y\in\sigma}
       \frac{|u(x)-u(y)|}{|x-y|_\sigma},
\tag{22.3}\label{eq:ch22-facewise-lipschitz}
\]

the maximum being over the finitely many faces of
\(\Delta_{\mathcal X}\).
\end{definition}

Arzel\`a--Ascoli requires equicontinuity on the compact space
\(\Delta_{\mathcal X}\), not merely on each of its faces separately.  The two
are equivalent here, and the elementary comparison is worth recording once,
because it is what converts the facewise estimate of
\Cref{thm:bfj-skeletal-equicontinuity} into a genuine modulus of continuity.

\begin{lemma}
\label{lem:ch22-path-metric}
Let \(\Delta\) be a connected finite polyhedral complex equipped with a
system of compatible norms, and define

\[
 d_\Delta(x,y):=\inf\Bigl\{\sum_{k=1}^{N}|x_{k}-x_{k-1}|:
     x_0=x,\ x_N=y,\ [x_{k-1},x_k]\subset\text{some face}\Bigr\}.
\tag{22.4}\label{eq:ch22-path-metric}
\]

where the length of a segment is measured in the norm of the smallest face
containing it.  Then \(d_\Delta\) is a metric inducing the topology of
\(\Delta\), and every \(u\) with
\(\operatorname{Lip}_{\mathrm f}(u)\leq C\) satisfies

\[
     |u(x)-u(y)|\leq C\,d_\Delta(x,y),
     \qquad x,y\in\Delta .
\tag{22.5}\label{eq:ch22-global-lipschitz}
\]

Moreover \(d_\Delta\) is bounded by a constant \(\delta(\Delta)\) depending
only on the complex and the chosen norms.
\end{lemma}

\begin{proof}
The infimum in \eqref{eq:ch22-path-metric} is over a nonempty set: any two
points of a connected finite complex are joined by a chain of segments each
contained in a face, since the complex is connected and every face is
convex, and it suffices to move along a chain of faces connecting the
carrier of \(x\) to the carrier of \(y\).  Symmetry and the triangle
inequality are clear by concatenation.  If \(d_\Delta(x,y)=0\) with
\(x\neq y\), choose a face \(\sigma\) containing \(x\) and a radius
\(r>0\) with \(|x-z|\geq r\) for all \(z\) in every face not containing
\(x\); such \(r\) exists because there are finitely many faces and the union
of those not containing \(x\) is compact and misses \(x\).  Any chain from
\(x\) to \(y\) either stays in the union of the faces containing \(x\), in
which case its length is at least the distance from \(x\) to \(y\) inside
the star of \(x\), or leaves that star and has length at least \(r\).  Both
alternatives contradict \(d_\Delta(x,y)=0\).  The same estimate shows that
\(d_\Delta\)-balls of small radius are neighborhoods and conversely, so the
metric induces the given topology.

For \eqref{eq:ch22-global-lipschitz}, apply the facewise bound to each
segment of a chain and sum; then take the infimum over chains.  Finally,
\(\Delta\) has finitely many faces, each of bounded diameter, and a fixed
chain of faces connects any two of them, so \(d_\Delta\) is bounded.
\end{proof}

\begin{remark}[Why the path metric and not a chordal one]
\label{rem:ch22-path-metric-choice}
A dual complex is usually pictured inside a Euclidean space, and one is
tempted to use the ambient distance.  The ambient distance is not intrinsic:
it depends on the embedding, whereas
\eqref{eq:ch22-path-metric} depends only on the integral affine structure
and the chosen compatible norms, hence, by
\Cref{prop:affine-structure-refinement}, is unchanged under refinement of
the complex.  Since our estimates must be compared across refinements, this
invariance is not a luxury.
\end{remark}

\subsection{Convexity on faces}

We recall from \Cref{chap:bfj-psh} the two structural facts that make the
estimates below meaningful, and record the elementary consequences used
later.

\begin{lemma}
\label{lem:ch22-convexity-consequences}
Let \(\varphi\in\PSH(X,\theta)\) and let \(\mathcal X\) be an SNC model on
which \(\theta\) is determined.  Then
\(u:=\varphi|_{\Delta_{\mathcal X}}\) is finite, continuous, and convex on
each face, and:
\begin{enumerate}
\item\label{itm:ch22-cvx:max} the maximum of \(u\) on a face is attained at
a vertex of that face, and the maximum of \(u\) on
\(\Delta_{\mathcal X}\) equals \(\sup_{X^{\An}}\varphi\);
\item\label{itm:ch22-cvx:slope} for each face \(\sigma\), each
\(v\in\sigma\) and each \(w\in V_\sigma\) with \(v+w\in\sigma\), the
one-sided derivative
\(\frac{\rd}{\rd t}\big|_{t=0+}u(v+tw)\) exists in
\([-\infty,+\infty)\) and is nondecreasing in \(t\) along the segment;
\item\label{itm:ch22-cvx:osc} if \(\max_{\Delta_{\mathcal X}}u=0\) and
\(\operatorname{Lip}_{\mathrm f}(u)\leq C\), then
\(-C\,\delta(\Delta_{\mathcal X})\leq u\leq0\).
\end{enumerate}
\end{lemma}

\begin{proof}
Finiteness, continuity and facewise convexity are
\Cref{thm:bfj-skeletal-characterization}\,(1).  Assertion
\cref{itm:ch22-cvx:max} is the standard fact that a convex function on a
simplex is bounded above by the affine interpolation of its vertex values,
together with \Cref{prop:ch21-sup-on-skeleton} for the second half.
Assertion \cref{itm:ch22-cvx:slope} is the monotonicity of difference
quotients of a one-variable convex function.  For
\cref{itm:ch22-cvx:osc}, pick \(x\) with \(u(x)=0\) and apply
\Cref{lem:ch22-path-metric}.
\end{proof}

\subsection{The simplex lemma}

Convexity bounds a function above by its vertex values, but it does not
bound it below, and it does not bound its slope near the boundary of a face.
The convex function \(t\mapsto-\sqrt t\) on \([0,1]\) is the reminder: it is
bounded, convex and continuous, has a finite value at both endpoints, and
has infinite one-sided derivative at \(t=0\).  Any estimate on the Lipschitz
constant of a convex function on a simplex must therefore feed on
information at the boundary.  The correct such information is the
derivative measured \emph{from} the opposite face \emph{toward} the vertex.

\begin{definition}[Derivative from the opposite face]
\label{def:ch22-directional-derivative}
Let \(\tau\) be a face with vertex \(e\), let \(\sigma\) be the opposite
face of \(\tau\), that is, the convex hull of the remaining vertices, and
let \(v\in\sigma\).  For a convex continuous \(u\) on \(\tau\) put

\[
 D_v u(e)
   :=\left.\frac{\rd}{\rd t}\right|_{t=0+}
      u\bigl((1-t)v+te\bigr)\in[-\infty,+\infty).
\tag{22.6}\label{eq:ch22-directional-derivative}
\]
\end{definition}

The parameter \(t\) in \eqref{eq:ch22-directional-derivative} is the affine
one, not arclength: \(D_vu(e)\) is the derivative of \(u\) at \(v\) in the
direction of the vector \(e-v\), whose length depends on \(v\).  This is the
normalization in which the estimate of \Cref{prop:ch22-slope-estimate} is an
intersection number, and we keep it.

\begin{definition}[Simplex gauge]
\label{def:ch22-simplex-gauge}
For a nondegenerate \(q\)-simplex \(\sigma\) with vertices
\(e_0,\ldots,e_q\) and direction space \(V\) equipped with a Euclidean norm,
put

\[
 W(\sigma):=\{e_j-e_i:\ i\neq j\},
 \qquad
 B(\sigma):=\operatorname{conv}W(\sigma),
\tag{22.7}\label{eq:ch22-gauge-body}
\]

and let

\[
 \Lambda(\sigma)
   :=\max_{|w|=1}\ \|w\|_{B(\sigma)}
   =\frac1{\operatorname{dist}\bigl(0,\partial B(\sigma)\bigr)},
\tag{22.8}\label{eq:ch22-gauge}
\]

where \(\|\cdot\|_{B}\) is the Minkowski gauge of the symmetric convex body
\(B\).
\end{definition}

The set \(W(\sigma)\) is symmetric and contains the basis
\(\{e_j-e_0\}_{j\geq1}\) of \(V\), so \(B(\sigma)\) is a symmetric convex
body and \eqref{eq:ch22-gauge} is a finite positive number.  The identity in
\eqref{eq:ch22-gauge} is the standard description of the gauge of a
symmetric body; we shall use it in the following form: for \(w\in V\),

\[
 \|w\|_{B(\sigma)}
 =\min\Bigl\{\sum_k\lambda_k:\ w=\sum_k\lambda_kw_k,\
       \lambda_k\geq0,\ w_k\in W(\sigma)\Bigr\}.
\tag{22.9}\label{eq:ch22-gauge-positive-combination}
\]

Indeed, if \(w=\sum\lambda_kw_k\) with \(\sum\lambda_k=t>0\), then
\(w/t\in B(\sigma)\), so \(\|w\|_{B(\sigma)}\leq t\); conversely, if
\(\|w\|_{B(\sigma)}=t\), then \(w/t\) is a convex combination of elements of
\(W(\sigma)\), and multiplying that combination by \(t\) exhibits \(w\) as
in \eqref{eq:ch22-gauge-positive-combination}.

\begin{lemma}
\label{lem:ch22-simplex}
Let \(u\) be convex and continuous on a nondegenerate \(q\)-simplex
\(\sigma\) with vertices \(e_0,\ldots,e_q\), and let \(\lambda\geq0\) be
such that

\[
 D_vu(e_i)\geq-\lambda
 \quad\text{for every vertex }e_i\text{ and every }
 v\in\RelInt(\sigma_i),
\tag{22.10}\label{eq:ch22-simplex-hypothesis}
\]

where \(\sigma_i\) is the face opposite \(e_i\).  Then:
\begin{enumerate}
\item\label{itm:ch22-simplex:vertices} the vertex values satisfy
\(|u(e_i)-u(e_j)|\leq\lambda\);
\item\label{itm:ch22-simplex:osc} the oscillation of \(u\) satisfies
\[
 \max_\sigma u-\min_\sigma u\leq(q+1)\lambda;
\tag{22.11}\label{eq:ch22-simplex-oscillation}
\]
\item\label{itm:ch22-simplex:lip} \(u\) is Lipschitz, with
\[
 |u(x)-u(y)|\leq(q+2)\,\Lambda(\sigma)\,\lambda\,|x-y|,
 \qquad x,y\in\sigma .
\tag{22.12}\label{eq:ch22-simplex-lipschitz}
\]
\end{enumerate}
No hypothesis other than \eqref{eq:ch22-simplex-hypothesis} is used; in
particular no a priori bound on \(u\) is assumed.
\end{lemma}

\begin{proof}
\emph{Step 1: the hypothesis on the closed opposite face.}
Fix \(i\) and \(v\in\RelInt(\sigma_i)\), and put
\(g(t)=u((1-t)v+te_i)\), a finite convex function on \([0,1]\).  Difference
quotients of a convex function decrease as the increment decreases, so

\[
 \frac{g(t)-g(0)}{t}\ \downarrow\ g'(0+)=D_vu(e_i)\geq-\lambda
 \qquad(t\downarrow0),
\]

whence

\[
 u\bigl((1-t)v+te_i\bigr)\geq u(v)-\lambda t,
 \qquad t\in[0,1].
\tag{22.13}\label{eq:ch22-segment-lower-bound}
\]

Both sides of \eqref{eq:ch22-segment-lower-bound} are continuous in \(v\)
because \(u\) is continuous on \(\sigma\), so
\eqref{eq:ch22-segment-lower-bound} holds for every \(v\) in the closed face
\(\sigma_i\); taking \(t\downarrow0\) again, the hypothesis
\eqref{eq:ch22-simplex-hypothesis} itself holds for every
\(v\in\sigma_i\), vertices included.

\emph{Step 2: vertex values and oscillation.}  Taking \(t=1\) in
\eqref{eq:ch22-segment-lower-bound} gives
\(u(e_i)\geq u(v)-\lambda\) for every \(v\in\sigma_i\), that is,

\[
 \max_{\sigma_i}u\leq u(e_i)+\lambda .
\tag{22.14}\label{eq:ch22-facet-max-bound}
\]

With \(v=e_j\) this is
\cref{itm:ch22-simplex:vertices}.  For the oscillation, we prove by
induction on \(q\) that
\(\min_\sigma u\geq\min_iu(e_i)-q\lambda\).  For \(q=0\) there is nothing to
prove.  For \(q\geq1\), every \(y\in\sigma\) lies on a segment
\([v,e_i]\) with \(v\in\sigma_i\) for any prescribed \(i\), so
\eqref{eq:ch22-segment-lower-bound} gives
\(u(y)\geq u(v)-\lambda\geq\min_{\sigma_i}u-\lambda\); the inductive
hypothesis applied to the \((q-1)\)-simplex \(\sigma_i\), whose vertices are
vertices of \(\sigma\), bounds \(\min_{\sigma_i}u\) below by
\(\min_iu(e_i)-(q-1)\lambda\).  Since
\(\max_\sigma u=\max_iu(e_i)\) by facewise convexity and
\(\max_iu(e_i)\leq\min_iu(e_i)+\lambda\) by
\cref{itm:ch22-simplex:vertices}, we obtain
\eqref{eq:ch22-simplex-oscillation}.

\emph{Step 3: two-sided bounds on subgradients.}  Let
\(x\in\RelInt(\sigma)\) and let \(p\) be a subgradient of \(u\) at \(x\),
so that \(u(y)\geq u(x)+\langle p,y-x\rangle\) for all \(y\in\sigma\); such
a \(p\) exists because \(u\) is convex and finite near \(x\) in the affine
hull of \(\sigma\).  Fix \(i\), let \(v\in\RelInt(\sigma_i)\) be the point
where the ray from \(e_i\) through \(x\) meets \(\sigma_i\), and write
\(x=(1-t_0)v+t_0e_i\) with \(t_0\in(0,1)\).  Restricting the subgradient
inequality to the segment shows that \(\langle p,e_i-v\rangle\) is a
subgradient at \(t_0\) of the convex function \(g\) of Step 1; hence

\[
 \langle p,e_i-v\rangle\geq g'(0+)=D_vu(e_i)\geq-\lambda .
\]

Since \(e_i-x=(1-t_0)(e_i-v)\) with \(1-t_0\in(0,1)\),

\[
 \langle p,e_i-x\rangle\geq-\lambda .
\tag{22.15}\label{eq:ch22-subgradient-lower}
\]

In the other direction, the subgradient inequality at \(y=e_i\) and
\eqref{eq:ch22-simplex-oscillation} give

\[
 \langle p,e_i-x\rangle\leq u(e_i)-u(x)\leq(q+1)\lambda .
\tag{22.16}\label{eq:ch22-subgradient-upper}
\]

\emph{Step 4: the gauge bound.}  For \(i\neq j\), subtracting
\eqref{eq:ch22-subgradient-lower} for the index \(i\) from
\eqref{eq:ch22-subgradient-upper} for the index \(j\) yields

\[
 \langle p,e_j-e_i\rangle
   =\langle p,e_j-x\rangle-\langle p,e_i-x\rangle
   \leq(q+1)\lambda+\lambda=(q+2)\lambda .
\tag{22.17}\label{eq:ch22-subgradient-edges}
\]

Thus \(\langle p,w\rangle\leq(q+2)\lambda\) for every \(w\in W(\sigma)\),
hence, by linearity, for every \(w\in B(\sigma)\), hence
\(\langle p,w\rangle\leq(q+2)\lambda\|w\|_{B(\sigma)}\) for every
\(w\in V\).  Taking for \(w\) the unit vector representing \(p\) under the
Euclidean structure gives
\(|p|\leq(q+2)\Lambda(\sigma)\lambda\).

\emph{Step 5: conclusion.}  A convex function whose subgradients at interior
points have norm at most \(M\) is \(M\)-Lipschitz on \(\RelInt(\sigma)\):
for interior \(x,y\) the subgradient inequalities at \(x\) and at \(y\) give
\(\langle p_y,x-y\rangle\leq u(x)-u(y)\leq\langle p_x,x-y\rangle\).  By
continuity of \(u\) on \(\sigma\) the bound extends to the closed simplex,
which is \eqref{eq:ch22-simplex-lipschitz}.
\end{proof}

\begin{remark}[How far the constant is from optimal]
\label{rem:ch22-simplex-constant}
The factor \(\Lambda(\sigma)\) in
\eqref{eq:ch22-simplex-lipschitz} cannot be improved: if \(p_0\) is a linear
form with \(\langle p_0,w\rangle\leq\lambda\) for all \(w\in W(\sigma)\) and
\(|p_0|=\lambda\Lambda(\sigma)\) --- such a form exists because the support
function of \(B(\sigma)\) attains the value \(\lambda\) in the direction
realizing the gauge in \eqref{eq:ch22-gauge} --- then the affine function
\(u=\langle p_0,\cdot-e_0\rangle\) satisfies
\eqref{eq:ch22-simplex-hypothesis} and has Lipschitz constant
\(\lambda\Lambda(\sigma)\).  The dimensional factor \(q+2\) is an artifact
of the proof.  In dimension one it can be removed entirely: on a segment,
\eqref{eq:ch22-simplex-hypothesis} says that the right derivative at one
endpoint is \(\geq-\lambda\) and the left derivative at the other is
\(\leq\lambda\), and convexity makes every intermediate slope lie between
these two, so that the sharp bound
\(\operatorname{Lip}(u)\leq\Lambda(\sigma)\lambda=\lambda/\ell\) holds.  We
use the sharp one-dimensional form in
\Cref{ex:ch22-modulus-segment} and the general form of
\Cref{lem:ch22-simplex} everywhere else; since only the existence of a
constant depending on \(\mathcal X\) is asserted in
\Cref{thm:bfj-skeletal-equicontinuity}, nothing later depends on the
difference.
\end{remark}

\begin{example}[The gauge of a segment and of the standard triangle]
\label{ex:ch22-gauge-computed}
Two computations are used verbatim in
\Cref{sec:ch22-worked-moduli}.

\emph{A segment.}  Let \(\sigma=[e_0,e_1]\) with
\(|e_1-e_0|=\ell\).  Then \(W(\sigma)=\{\pm(e_1-e_0)\}\), so
\(B(\sigma)\) is the segment of half-length \(\ell\) and

\[
 \Lambda(\sigma)=\frac1\ell .
\tag{22.18}\label{eq:ch22-gauge-segment}
\]

In the coordinate \(t\in[0,1]\) of \Cref{ex:ch22-edge-normalization} one has
\(\ell=1\) and \(\Lambda=1\): a convex function on the unit edge whose
endpoint derivatives toward the opposite endpoint are \(\geq-\lambda\) has
all slopes in \([-\lambda,\lambda]\).

\emph{The standard triangle.}  Let
\(\sigma=\operatorname{conv}\{(0,0),(1,0),(0,1)\}\subset\mathbf R^2\) with
the Euclidean norm; by \Cref{ex:dual-complex-triangle} and
\Cref{ex:toric-triangle-chart} this is the compatible normalization of the
dual complex of the chart \(xyz=\pi\), in the coordinates
\((\alpha_x,\alpha_y)\).  Then

\[
 W(\sigma)=\{\pm(1,0),\ \pm(0,1),\ \pm(1,-1)\},
\]

and \(B(\sigma)\) is the hexagon whose vertices, in cyclic order, are
\[
 (1,0),\ (0,1),\ (-1,1),\ (-1,0),\ (0,-1),\ (1,-1).
\]
Its boundary meets the
diagonal direction \((1,1)/\sqrt2\) on the edge joining \((1,0)\) to
\((0,1)\), at distance \(1/\sqrt2\) from the origin, while every other edge
lies at distance \(1\).  Hence

\[
 \Lambda(\sigma)=\sqrt2 .
\tag{22.19}\label{eq:ch22-gauge-triangle}
\]
\end{example}

\section{The vertex estimate}
\label{sec:ch22-vertex}

Throughout this section and the next, \(\mathcal X\) is an SNC model on which
\(\theta\) is determined, with

\[
            \mathcal X_s=\sum_{i\in I}b_iE_i,
            \qquad e_i=x_{E_i}\in\Delta_{\mathcal X},
\tag{22.20}\label{eq:ch22-vertices}
\]

where \(e_i=b_i^{-1}\ord_{E_i}\) is the divisorial point of \(E_i\), as in
\Cref{rem:divisorial-normalization}.  We fix once and for all a relatively
ample \emph{line bundle} \(\mathcal A\) on \(\mathcal X\); integrality of
\(\mathcal A\) is what makes the intersection numbers below integers, and it
costs nothing since a positive multiple of a relatively ample
\(\mathbf Q\)-class is such a bundle.

\subsection{The incidence data}

\begin{lemma}
\label{lem:ch22-incidence-signs}
Set

\[
 w_{ij}:=b_j\bigl(E_i\cdot E_j\cdot\mathcal A^{n-1}\bigr),
 \qquad i,j\in I .
\tag{22.21}\label{eq:ch22-incidence-weights}
\]

Then:
\begin{enumerate}
\item\label{itm:ch22-inc:integrality} every \(w_{ij}\) is an integer;
\item\label{itm:ch22-inc:offdiag}
\[
 w_{ij}\geq0\quad(i\neq j),
\tag{22.22}\label{eq:ch22-off-diagonal-positive}
\]
with \(w_{ij}\geq1\) whenever \(E_i\cap E_j\neq\varnothing\);
\item\label{itm:ch22-inc:rowsum}
\[
 \sum_{j\in I}w_{ij}
       =\bigl(E_i\cdot\mathcal X_s\cdot\mathcal A^{n-1}\bigr)=0;
\tag{22.23}\label{eq:ch22-row-sum-zero}
\]
\item\label{itm:ch22-inc:connected} the graph on \(I\) with an edge
\(\{i,j\}\) whenever \(E_i\cap E_j\neq\varnothing\) is connected.
\end{enumerate}
\end{lemma}

\begin{proof}
Integrality holds because \(\mathcal A\) and the \(E_i\) are integral
divisors on the projective \(R\)-scheme \(\mathcal X\) and the special fibre
is proper over the residue field.  For
\cref{itm:ch22-inc:offdiag}, distinct components of an SNC divisor meet
transversally along an effective cycle \(E_i\cap E_j\) of dimension
\(n-1\), and \(\mathcal A\) restricts to an ample bundle on each of its
components; hence the intersection number is a sum of positive degrees, and
it is a positive integer as soon as the cycle is nonzero.  Assertion
\cref{itm:ch22-inc:rowsum} holds because
\(\mathcal O_{\mathcal X}(\mathcal X_s)\) is trivial: the special fibre is
the divisor of \(\pi\).  This is the same row-sum identity that produced the
graph Laplacian in \eqref{eq:ch20-graph-laplacian} and the relative Hodge
form \eqref{eq:ch20-hodge-laplacian}.

For \cref{itm:ch22-inc:connected}, recall from the proof of
\Cref{thm:ch20-ddc-exact} that \(\mathcal X_s\) is connected: otherwise the
theorem on formal functions would split
\(H^0(\mathcal X,\mathcal O_{\mathcal X})\) into a nontrivial product, while
its restriction to the generic fibre is a domain because \(X\) is connected
and smooth.  A connected SNC divisor has connected incidence graph, since a
connected topological space that is a finite union of closed connected
subsets has connected nerve.
\end{proof}

The three signs in \Cref{lem:ch22-incidence-signs} are the whole mechanism of
this section: a nonnegative off-diagonal part, a vanishing row sum, and
connectedness.  The matrix \((w_{ij})\) is, up to the positive factors
\(b_i\), the Laplacian of a connected weighted graph, and everything below is
a maximum-principle argument for that Laplacian in the presence of the
inhomogeneous term supplied by \(\theta\).

\subsection{The estimate}

\begin{proposition}
\label{prop:ch22-vertex-estimate}
Put

\[
 c:=\max_{i\in I}\bigl|\bigl(\theta_{\mathcal X}\cdot E_i\cdot
       \mathcal A^{n-1}\bigr)\bigr|,
 \qquad
 \lambda:=\max_{i\in I}\Bigl(-b_i\bigl(E_i^2\cdot
       \mathcal A^{n-1}\bigr)\Bigr).
\tag{22.24}\label{eq:ch22-lambda}
\]

If \(\#I\geq2\), then \(\lambda\geq1\), and every \(\theta\)-psh model
function \(\varphi\) normalized by \eqref{eq:ch22-normalization} satisfies

\[
 0\geq\varphi(e_i)\geq-C_{\mathrm{vert}},
 \qquad
 C_{\mathrm{vert}}:=c\sum_{m=0}^{\#I-2}\lambda^m,
 \qquad i\in I .
\tag{22.25}\label{eq:ch22-vertex-lower-bound}
\]

If \(\#I=1\), then \(\Delta_{\mathcal X}\) is a point and
\(\varphi(e_i)=0\).
\end{proposition}

\begin{proof}
The case \(\#I=1\) is immediate from
\eqref{eq:ch22-normalization} and
\Cref{prop:ch21-sup-on-skeleton}, so assume \(\#I\geq2\).  By
\eqref{eq:ch22-row-sum-zero},

\[
 -b_i\bigl(E_i^2\cdot\mathcal A^{n-1}\bigr)
 =-w_{ii}
 =\sum_{j\neq i}w_{ij}
 =\sum_{j\neq i}b_j\bigl(E_i\cdot E_j\cdot\mathcal A^{n-1}\bigr),
\]

a sum of nonnegative integers which is positive as soon as \(E_i\) meets
another component.  Connectedness and \(\#I\geq2\) guarantee that some \(E_i\) does,
so \(\lambda\geq1\).

\paragraph{The vertex calculation.}
Let \(\varphi\) be a \(\theta\)-psh model function.  Choose a domination
\(\pi:\mathcal Y\to\mathcal X\) and a vertical \(\mathbf R\)-Cartier divisor
\(G\) on \(\mathcal Y\) with \(\varphi=\varphi_G\) and
\(\pi^*\theta_{\mathcal X}+[G]\) relatively nef; such data exist by
\Cref{cor:ch20-reference-change} and the definition of model
semipositivity.  Adding to \(G\) a real multiple of \(\mathcal Y_s\) changes
\(\varphi\) by a constant and changes no positivity, so we may and do assume
the normalization \eqref{eq:ch22-normalization}.  By
\Cref{lem:ch21-model-skeletal-properties} the restriction
\(\varphi|_{\Delta_{\mathcal X}}\) is convex on each face, and by
\Cref{prop:ch21-sup-on-skeleton} its maximum is the global supremum; a
convex function on a simplex attains its maximum at a vertex, so

\[
                         \max_{i\in I}\varphi(e_i)=0 .
\tag{22.26}\label{eq:ch22-vertex-max}
\]

The coefficient of \(E_j\) in \(\pi_*G\) is
\(\ord_{E_j}(G)=b_j\varphi_G(e_j)\), because
\(e_j=b_j^{-1}\ord_{E_j}\) and \(\varphi_G(x)=v_x(f_G)\) by
\eqref{eq:ch20-sign}.  Hence

\[
             \pi_*G=\sum_{j\in I}b_j\varphi(e_j)E_j.
\tag{22.27}\label{eq:ch22-pushforward-divisor}
\]

The class \(\pi^*\theta_{\mathcal X}+[G]\) is nef on the special fibre.
Intersect it with the effective divisor \(\pi^*E_i\) and with
\((\pi^*\mathcal A)^{n-1}\).  The result is nonnegative: restricting to the
components of the effective cycle \(\pi^*E_i\), one is intersecting a nef
class with \(n-1\) nef classes on a projective variety over the residue
field.  Now apply the projection formula
\eqref{eq:appF-projection} to push the computation down to \(\mathcal X\),
using that \(\theta_{\mathcal X}\), \(E_i\) and \(\mathcal A\) are pulled
back and that \(\pi_*\) of \(G\) is \eqref{eq:ch22-pushforward-divisor}.
The result is

\[
 \sum_{j\in I} b_j\varphi(e_j)
       \bigl(E_i\cdot E_j\cdot\mathcal A^{n-1}\bigr)
 \geq
 -\bigl(\theta_{\mathcal X}\cdot E_i\cdot\mathcal A^{n-1}\bigr)
 \geq-c,
\tag{22.28}\label{eq:ch22-vertex-intersection}
\]

valid for every \(i\in I\).  All quantities other than the vertex values
\(\varphi(e_j)\) depend only on the fixed data
\((\mathcal X,\theta,\mathcal A)\).

\paragraph{Propagation along a path.}
By \eqref{eq:ch22-vertex-max} all vertex values are nonpositive and at least
one vanishes.  Choose a vertex \(i_0\) with \(\varphi(e_{i_0})=0\) and a
vertex \(i_M\) with \(\varphi(e_{i_M})=\min_i\varphi(e_i)\), and, using
\Cref{lem:ch22-incidence-signs}\,\cref{itm:ch22-inc:connected}, a path

\[
             i_0,i_1,\ldots,i_M
             \quad\text{with}\quad
             E_{i_m}\cap E_{i_{m+1}}\neq\varnothing
\tag{22.29}\label{eq:ch22-vertex-path}
\]

in the incidence graph with pairwise distinct vertices; then
\(M\leq\#I-1\).  Fix \(m<M\) and apply
\eqref{eq:ch22-vertex-intersection} at \(i=i_m\).  Separating the diagonal
term and using \eqref{eq:ch22-row-sum-zero},

\[
 -\Bigl(\sum_{j\neq i_m}w_{i_mj}\Bigr)\varphi(e_{i_m})
 +\sum_{j\neq i_m}w_{i_mj}\varphi(e_j)\geq-c .
\]

Every term of the second sum is nonpositive, since \(w_{i_mj}\geq0\) and
\(\varphi\leq0\); discarding all of them except the one with
\(j=i_{m+1}\) only increases the right-hand side of the rearranged
inequality, so

\[
 \Bigl(\sum_{j\neq i_m}w_{i_mj}\Bigr)\varphi(e_{i_m})
 \leq c+w_{i_mi_{m+1}}\varphi(e_{i_{m+1}})
 \leq c+\varphi(e_{i_{m+1}}),
\]

the last step because \(w_{i_mi_{m+1}}\geq1\) and
\(\varphi(e_{i_{m+1}})\leq0\).  Since
\(1\leq\sum_{j\neq i_m}w_{i_mj}\leq\lambda\) and
\(\varphi(e_{i_m})\leq0\), the left-hand side is at least
\(\lambda\varphi(e_{i_m})\), whence

\[
               \lambda\,\varphi(e_{i_m})
               \leq c+\varphi(e_{i_{m+1}}).
\tag{22.30}\label{eq:ch22-path-recursion}
\]

Equivalently \(\varphi(e_{i_{m+1}})\geq\lambda\varphi(e_{i_m})-c\).
Starting from \(\varphi(e_{i_0})=0\) and iterating \(M\leq\#I-1\) times,

\[
 \varphi(e_{i_M})
 \geq-c\sum_{m=0}^{M-1}\lambda^m
 \geq-c\sum_{m=0}^{\#I-2}\lambda^m
 =-C_{\mathrm{vert}} ,
\]

which is \eqref{eq:ch22-vertex-lower-bound}, since \(e_{i_M}\) realizes the
minimum.
\end{proof}

\begin{remark}[What the vertex estimate is, and is not]
\label{rem:ch22-vertex-estimate-nature}
Three features of the proof deserve to be recorded, because they are what
make the estimate usable later.

First, no compactness theorem has been used: the bound is a finite
computation with the fixed matrix \((w_{ij})\) and the fixed vector
\(\bigl((\theta_{\mathcal X}\cdot E_i\cdot\mathcal A^{n-1})\bigr)_i\).  This
is why the estimate may be quoted inside the proof of
\Cref{thm:ch21-decreasing-stability} without circularity.

Second, the constant is uniform over \emph{all} \(\theta\)-psh model
functions, however high the model on which they are determined.  The reason
is that a determination on a higher model enters only through the
pushforward \eqref{eq:ch22-pushforward-divisor}, whose coefficients are the
values on the fixed complex \(\Delta_{\mathcal X}\).

Third, the geometric input is exactly the pair of sign conditions
\eqref{eq:ch22-off-diagonal-positive}--\eqref{eq:ch22-row-sum-zero} and the
connectedness of the special fibre.  In particular, the estimate is
insensitive to the residue characteristic; see
\Cref{prop:ch22-characteristic-free-estimates}.
\end{remark}

\begin{remark}[The exponential constant is an artifact]
\label{rem:ch22-exponential-constant}
The constant \(C_{\mathrm{vert}}\) grows exponentially in the number of
components.  This is not an accident of a wasteful step: a chain of
components along which the reference degrees vanish transmits a slope
without damping, and the recursion \eqref{eq:ch22-path-recursion} with
\(\lambda>1\) then amplifies.  For a \emph{semistable} model in the strict
sense one has \(b_i=1\), but in relative dimension greater than one the
weights
\(w_{ij}=(E_i\cdot E_j\cdot\mathcal A^{n-1})\) still depend on the chosen
polarization and need not lie in \(\{0,1\}\).  Thus
\(\lambda=\max_i\sum_{j\ne i}w_{ij}\) is the maximum weighted valence in
general.  Unit weights occur for individual nodes of a semistable curve; if
parallel nodes are combined into one adjacency, their contributions add.
The estimate is a weighted discrete maximum principle with the same
exponential shape.  For the two-component curve of
\Cref{ex:ch22-two-vertex} the constant is \(C_{\mathrm{vert}}=
\max\{d_0,d_1\}\), and \eqref{eq:ch22-segment-sharp-constants} exhibits a
normalized \(\theta\)-psh function attaining it; so the exponential shape,
not the constant itself, is what is potentially lossy.
\end{remark}

\section{Toroidal techniques and nef approximation}
\label{sec:ch22-toroidal}

The vertex estimate used only the incidence numbers of a fixed model.  A
bound on slopes cannot: a slope is a derivative at a point that is not a
vertex, and the model must be refined before that derivative becomes an
intersection number.  The refinements needed are toroidal, they are
prescribed by convex piecewise affine data on the dual complex, and the
positivity they must transport is relative nefness.  This section assembles
the four tools that make this transport legitimate.  All of them are
independent of the residue characteristic, a point made precise in
\Cref{prop:ch22-characteristic-free-estimates}.

\subsection{Approximation of relatively nef divisors by relatively ample ones}

\begin{lemma}
\label{lem:ch22-nef-approximation}
Let \(\mu:\mathcal Y\to\mathcal Z\) be a projective birational morphism of
normal projective \(R\)-models of \(X\), and let \(G\) be a vertical
\(\mathbf R\)-Cartier divisor on \(\mathcal Y\) that is \(\mu\)-nef.  For
every \(\mu\)-ample \(\mathbf R\)-Cartier divisor \(A\) and every
\(\varepsilon>0\), the divisor \(G+\varepsilon A\) is \(\mu\)-ample.
Consequently \(G\) is the limit, in the finite-dimensional real vector space
of vertical \(\mathbf R\)-Cartier divisor classes modulo relative numerical
equivalence, of \(\mu\)-ample \(\mathbf Q\)-Cartier divisors, and for each
such divisor \(A'\) there is \(m\geq1\) with
\(\mathcal O_{\mathcal Y}(mA')\) generated relative to \(\mathcal Z\).
\end{lemma}

\begin{proof}
Relative nefness and relative ampleness are tested on curves contracted by
\(\mu\); by the relative Kleiman criterion \cite[Section~1.4]{KM08}, the
relative nef cone is the closure of the relative ample cone in the
finite-dimensional space \(N^1(\mathcal Y/\mathcal Z)_{\mathbf R}\), and the
ample cone is open.  Hence the sum of a nef class and a positive multiple of
an ample one is ample, which is the first assertion, and \(G\) is a limit of
ample rational classes because the rational points of an open cone are dense
in it.  For the last statement, relative ampleness of \(A'\) gives, after
clearing denominators, a line bundle \(\mathcal O_{\mathcal Y}(mA')\) that is
relatively very ample for \(m\) divisible enough, hence generated by its
relative sections.
\end{proof}

\begin{remark}[Where the finite dimension comes from]
The vector space in \Cref{lem:ch22-nef-approximation} is finite dimensional
because \(\mathcal Y_s\) has finitely many components and relative numerical
classes of vertical divisors are detected by finitely many curve classes;
this is the same finiteness that made \(N^1(\mathcal X/R)\) finite
dimensional in \Cref{chap:bfj-model-forms}.  Nothing here requires the
existence of an SNC model.
\end{remark}

\subsection{The negativity lemma}

\begin{lemma}
\label{lem:ch22-negativity}
Let \(\mu:\mathcal Y\to\mathcal Z\) be a projective birational morphism of
normal projective \(R\)-models of \(X\), with \(\mathcal Z\) vertically
\(\mathbf Q\)-factorial, and let \(G\) be a vertical \(\mathbf R\)-Cartier
divisor on \(\mathcal Y\) that is \(\mu\)-nef.  Then

\[
                 \mu^*\mu_*G-G\geq0 .
\tag{22.31}\label{eq:ch22-negativity}
\]
\end{lemma}

\begin{proof}
Both sides of \eqref{eq:ch22-negativity} are continuous and linear in \(G\)
on the finite-dimensional space of vertical \(\mathbf R\)-divisors, and the
conclusion is a closed condition; by
\Cref{lem:ch22-nef-approximation} it therefore suffices to treat a
\(\mu\)-ample \(\mathbf Q\)-Cartier divisor \(G\).  Rescaling, assume that
\(\mathcal O_{\mathcal Y}(mG)\) is a line bundle generated relative to
\(\mathcal Z\), and set
\(\mathfrak a:=\mu_*\mathcal O_{\mathcal Y}(mG)\), a coherent fractional
ideal on \(\mathcal Z\).

Relative global generation says exactly that

\[
       \mathfrak a\,\mathcal O_{\mathcal Y}
       =\mathcal O_{\mathcal Y}(mG).
\tag{22.32}\label{eq:ch22-generation-identity}
\]

On the other hand, a section of \(\mathcal O_{\mathcal Y}(mG)\) over
\(\mu^{-1}(U)\) is a rational function \(f\) with
\(\divis_{\mathcal Y}(f)+mG\geq0\); pushing forward divisors preserves
effectivity, so \(\divis_{\mathcal Z}(f)+m\mu_*G\geq0\), that is,

\[
 \mathfrak a\subset\mathcal O_{\mathcal Z}(m\mu_*G).
\tag{22.33}\label{eq:ch22-pushforward-inclusion}
\]

Because \(\mathcal Z\) is vertically \(\mathbf Q\)-factorial, \(\mu_*G\) is
\(\mathbf Q\)-Cartier and \(\mu^*\mu_*G\) is defined; pulling back
\eqref{eq:ch22-pushforward-inclusion} and using
\eqref{eq:ch22-generation-identity} gives
\(\mathcal O_{\mathcal Y}(mG)\subset\mathcal O_{\mathcal Y}(m\mu^*\mu_*G)\),
that is, \(mG\leq m\mu^*\mu_*G\).  Dividing by \(m\) proves
\eqref{eq:ch22-negativity}.  This is \cite[Lemma~1.6]{BFJ16}.
\end{proof}

\begin{warning}{\(\mathbf Q\)-factoriality is a hypothesis, not a formality}
Without vertical \(\mathbf Q\)-factoriality of \(\mathcal Z\) the symbol
\(\mu^*\mu_*G\) has no meaning, and the conclusion of
\Cref{lem:ch22-negativity} cannot even be stated.  This is why the slope
estimate is run on the model produced by
\Cref{thm:skeleton-subdivision}, which is normal and vertically
\(\mathbf Q\)-factorial by construction, rather than on an arbitrary
refinement.  The order relation in \eqref{eq:ch22-negativity} is
coefficientwise effectivity of divisors, not pseudo-effectivity of numerical
classes; the two orders are separated in
\Cref{app:intersection-bdivisors}.
\end{warning}

\subsection{Shrinking a star}

We now construct the subdivision used by the slope estimate.  Let \(\sigma\)
be a face of \(\Delta_{\mathcal X}\) with vertex set \(J\subset I\), let

\[
 L:=\{j\in I:\ \sigma\cup\{e_j\}\ \text{spans a face of}\
       \Delta_{\mathcal X}\}\supset J
\tag{22.34}\label{eq:ch22-star-index}
\]

index the vertices of the star of \(\sigma\), and let \(v\in\RelInt(\sigma)\)
be a rational point, written in barycentric coordinates as

\[
 v=\sum_{j\in J}s_je_j,
 \qquad s_j>0,\quad\sum_{j\in J}s_j=1 .
\tag{22.35}\label{eq:ch22-barycentric-v}
\]

Here, as in \Cref{ex:dual-complex-triangle}, the barycentric coordinate
\(t_i\) of a point of \(\Delta_{\mathcal X}\) is
\(t_i=b_i\alpha_i\), so that \(t_i\) is the restriction to
\(\Delta_{\mathcal X}\) of the model function \(\chi_i:=\varphi_{b_iE_i}\)
of \Cref{prop:vertical-divisor-affine-iso}; in particular each \(t_i\) is
nonnegative, affine on every face, and \(\chi_i(e_j)=\delta_{ij}\).

\begin{lemma}
\label{lem:ch22-star-shrinking}
For every rational \(\varepsilon\in(0,1)\) put

\[
                 e'_j:=\varepsilon e_j+(1-\varepsilon)v,
                 \qquad j\in L .
\tag{22.36}\label{eq:ch22-star-shrink}
\]

There is a rational projective simplicial subdivision \(\Delta'\) of
\(\Delta_{\mathcal X}\) with the following properties.
\begin{enumerate}
\item\label{itm:ch22-star:vertices} The points \(e'_j\), \(j\in L\), are
vertices of \(\Delta'\).
\item\label{itm:ch22-star:face} The simplex
\(\sigma'_J:=\operatorname{conv}\{e'_j:j\in J\}\) is a face of \(\Delta'\),
it is contained in \(\sigma\) and has the same dimension, and
\(v\in\RelInt(\sigma'_J)\).
\item\label{itm:ch22-star:star} The faces of \(\Delta'\) containing
\(\sigma'_J\) are exactly the simplices
\(\operatorname{conv}\{e'_j:j\in J\cup M\}\), where \(\sigma\cup\{e_j\}_{j\in
M}\) spans a face of \(\Delta_{\mathcal X}\); in particular the vertices of
\(\Delta'\) adjacent to \(\sigma'_J\) are exactly the \(e'_j\) with
\(j\in L\).
\end{enumerate}
Consequently, by \Cref{thm:skeleton-subdivision}, there is a vertical blowup
\(\rho:\mathcal X'\to\mathcal X\) with \(\mathcal X'\) normal and vertically
\(\mathbf Q\)-factorial whose strata are indexed by the faces of
\(\Delta'\), and on which vertical Cartier divisors restrict to affine
functions on the faces of \(\Delta'\).
\end{lemma}

\begin{proof}
Consider the function

\[
 h:=\max\Bigl\{\max_{j\in J}\Bigl(-\frac{\chi_j}{s_j}\Bigr),
       \ -(1-\varepsilon)\Bigr\}
\tag{22.37}\label{eq:ch22-support-function}
\]

on \(\Delta_{\mathcal X}\).  Each \(\chi_j\) is a rational model function,
affine on every face; hence \(h\) is a rational piecewise affine function,
convex on every face as a maximum of affine functions, and its domains of
affineness define a subdivision \(\Delta_h\) of
\(\Delta_{\mathcal X}\), projective by construction.

Fix a maximal face \(F\) of \(\Delta_{\mathcal X}\) containing \(\sigma\),
with vertex set \(J\cup M\).  On \(F\) the region where the constant branch
realizes the maximum is

\[
 S_F:=\{t\in F:\ t_j\geq(1-\varepsilon)s_j\ \text{for all }j\in J\}.
\]

Since \(\sum_{i\in J\cup M}t_i=1\) and
\(\sum_{j\in J}(1-\varepsilon)s_j=1-\varepsilon\), the region \(S_F\) is the
image of \(F\) under the homothety of ratio \(\varepsilon\) centered at
\(v\); its vertices are precisely the points \(e'_j\), \(j\in J\cup M\), and
it is a simplex.  Indeed the barycentric coordinates of \(e'_k\) are

\[
              t_j(e'_k)=\chi_j(e'_k)
              =\varepsilon\delta_{jk}+(1-\varepsilon)s_j,
              \qquad j\in J\cup M,\ k\in J\cup M,
\tag{22.38}\label{eq:ch22-coordinate-values}
\]

with \(s_j:=0\) for \(j\notin J\), so \(e'_k\in S_F\) for every \(k\), and a
point of \(F\) lies in \(S_F\) if and only if it is a convex combination of
the \(e'_k\).  The faces of \(S_F\) include
\(\sigma'_J=\operatorname{conv}\{e'_j:j\in J\}\), which is therefore a common
face of the cells \(S_F\) for all maximal \(F\supset\sigma\), and no cell of
\(\Delta_h\) outside the shrunk star contains \(\sigma'_J\), because such a
cell lies in a region where some \(-\chi_j/s_j\) exceeds
\(-(1-\varepsilon)\) while every point of \(\sigma'_J\) satisfies
\(\chi_j\geq(1-\varepsilon)s_j\), with equality only on the boundary faces of
\(S_F\) not containing \(v\).  This proves
\cref{itm:ch22-star:vertices}--\cref{itm:ch22-star:star} for \(\Delta_h\),
except that the cells of \(\Delta_h\) outside the shrunk star need not be
simplices, and \(\sigma'_J\subset\sigma\) with
\(\dim\sigma'_J=\dim\sigma\) and \(v=\sum_{j\in J}s_je'_j\in
\RelInt(\sigma'_J)\) by \eqref{eq:ch22-star-shrink} and
\eqref{eq:ch22-barycentric-v}.

To make the subdivision simplicial, refine \(\Delta_h\) by a pulling
triangulation at its vertices: fix an ordering of the finitely many vertices
of \(\Delta_h\) and add to \(h\) a sufficiently small generic rational convex
piecewise affine function with the same vertex set.  The resulting function
is again convex and rational, and its domains of affineness form a
simplicial subdivision \(\Delta'\) refining \(\Delta_h\) and using no new
vertices.  A cell of \(\Delta_h\) that is already a simplex is not
subdivided by such a refinement, so the shrunk-star cells \(S_F\), and with
them \(\sigma'_J\) and the incidence description
\cref{itm:ch22-star:star}, survive unchanged.  Finally
\(\Delta'\) is projective, being defined by a convex rational piecewise
affine function, and \Cref{thm:skeleton-subdivision} applies.
\end{proof}

\begin{remark}[What the shrinking is for]
\label{rem:ch22-why-shrink}
The point of \eqref{eq:ch22-star-shrink} is to convert a directional
derivative at \(v\) into a divisor coefficient.  After shrinking, the value
of a psh function at the new vertex \(e'_j\) is
\(\varphi(v)+\varepsilon D_v\varphi(e_j)\) whenever \(\varphi\) is affine on
the short segment \([v,e'_j]\); so the vector of derivatives
\((D_v\varphi(e_j))_{j\in L}\) is, up to the factor \(\varepsilon\) and the
constant \(\varphi(v)\), the vector of vertex values on the new model.  The
vertex mechanism of \Cref{sec:ch22-vertex} can then be applied on
\(\mathcal X'\), where it becomes an inequality between intersection numbers
on the stratum \(E'_J\).  The price is that \(\mathcal X'\) is only
\(\mathbf Q\)-factorial and toroidal, not SNC; this is precisely why
\Cref{lem:ch22-negativity} was stated in that generality.
\end{remark}

\subsection{Two positivity transfers}

\begin{lemma}
\label{lem:ch22-pullback-positivity}
Let \(g:W\to V\) be a surjective morphism of projective varieties over a
field, let \(\alpha\) be a numerical class on \(V\), and suppose that

\[
                 g^*\alpha=\gamma+F
\]

with \(\gamma\) nef and \(F\) an effective \(\mathbf R\)-divisor class on
\(W\).  Then, for every nef class \(\beta\) on \(V\),

\[
                      \bigl(\alpha\cdot\beta^{\dim V-1}\bigr)_V\geq0 .
\tag{22.39}\label{eq:ch22-pullback-pseudoeffective}
\]
\end{lemma}

\begin{proof}
Both sides are unchanged by base change to an algebraic closure, so assume
the ground field algebraically closed and infinite.  Set \(d=\dim V\).  Cut
\(W\) by \(\dim W-d\) general members of a very ample linear system: the
resulting subvariety \(W'\) has dimension \(d\), maps onto \(V\), and is not
contained in the support of \(F\), because general members avoid any
prescribed proper closed subset in the relevant sense and the fibres of
\(g\) are cut down to finite sets.  Thus \(g|_{W'}:W'\to V\) is generically
finite and surjective, of some degree \(\delta>0\), and \(F|_{W'}\) is an
effective class.  The projection formula
\eqref{eq:appF-projection} gives

\[
 \begin{aligned}
 \delta\bigl(\alpha\cdot\beta^{d-1}\bigr)_V
 &=\Bigl((g|_{W'})^*\alpha\cdot
    \bigl((g|_{W'})^*\beta\bigr)^{d-1}\Bigr)_{W'}\\
 &=\Bigl(\gamma|_{W'}\cdot(g^*\beta)^{d-1}\Bigr)\\
 &\quad+\Bigl(F|_{W'}\cdot(g^*\beta)^{d-1}\Bigr).
 \end{aligned}
\]

The first term is nonnegative because \(\gamma\) and \(g^*\beta\) are nef;
the second is nonnegative because \(F|_{W'}\) is effective and
\(g^*\beta\) is nef.  Divide by \(\delta\).
\end{proof}

\begin{lemma}
\label{lem:ch22-restriction-identity}
Keep the notation of \Cref{lem:ch22-star-shrinking}, let
\(\rho:\mathcal X'\to\mathcal X\) realize \(\Delta'\), let
\(\mu:\mathcal Y\to\mathcal X'\) dominate a determination \(G\) of a model
function \(\varphi\), and let \(E'_J\) be the stratum of \(\mathcal X'\)
associated with \(\sigma'_J\).  Assume that \(\varphi\) is affine on
\(\sigma'_J\) and on each segment \([v,e'_j]\), \(j\in L\).  Then

\[
 \left.
 \rho^*\!\left(\sum_{j\in L}D_v\varphi(e_j)\,b_jE_j\right)
 \right|_{E'_J}
       =\left.\mu_*G\right|_{E'_J}.
\tag{22.40}\label{eq:ch22-slope-divisor-identity}
\]
\end{lemma}

\begin{proof}
Write \(b'_j\) for the multiplicity of the component of
\(\mathcal X'_s\) attached to \(e'_j\).  A vertical divisor on
\(\mathcal X'\) with model function \(\psi\) has coefficient
\(b'_j\psi(e'_j)\) along that component, by
\eqref{eq:ch20-sign}.  Applying this to
\(\rho^*(b_iE_i)\), whose model function is \(\chi_i\), and restricting to
\(E'_J\), only the components adjacent to \(E'_J\) survive; these are indexed
by \(L\) by \Cref{lem:ch22-star-shrinking}\,\cref{itm:ch22-star:star}.  Hence,
by \eqref{eq:ch22-coordinate-values},

\[
 \rho^*(b_iE_i)|_{E'_J}
   =\sum_{j\in L}\bigl(\varepsilon\delta_{ij}
        +(1-\varepsilon)s_i\bigr)b'_jE'_j\big|_{E'_J}.
\tag{22.41}\label{eq:ch22-pullback-restriction}
\]

Since \(\varphi\) is affine on \([v,e'_i]\) and
\(e'_i=(1-\varepsilon)v+\varepsilon e_i\),

\[
 D_v\varphi(e_i)
   =\varepsilon^{-1}\bigl(\varphi(e'_i)-\varphi(v)\bigr),
       \qquad i\in L .
\tag{22.42}\label{eq:ch22-derivative-as-value}
\]

Insert \eqref{eq:ch22-pullback-restriction} and
\eqref{eq:ch22-derivative-as-value} into the left-hand side of
\eqref{eq:ch22-slope-divisor-identity}.  The coefficient of
\(b'_jE'_j|_{E'_J}\) becomes

\[
 \varepsilon D_v\varphi(e_j)
 +(1-\varepsilon)\sum_{i\in L}s_iD_v\varphi(e_i)
 =\bigl(\varphi(e'_j)-\varphi(v)\bigr)
 +\frac{1-\varepsilon}{\varepsilon}
   \Bigl(\sum_{i\in J}s_i\varphi(e'_i)-\varphi(v)\Bigr),
\]

because \(s_i=0\) for \(i\notin J\).  The last bracket vanishes: by
\eqref{eq:ch22-star-shrink} and \eqref{eq:ch22-barycentric-v} one has
\(v=\sum_{i\in J}s_ie'_i\) with \(\sum_{i\in J}s_i=1\), and \(\varphi\) is
affine on \(\sigma'_J\).  Therefore

\[
 \left.
 \rho^*\!\left(\sum_{j\in L}D_v\varphi(e_j)b_jE_j\right)\right|_{E'_J}
 =\sum_{j\in L}\bigl(\varphi(e'_j)-\varphi(v)\bigr)b'_jE'_j\big|_{E'_J}
 =\bigl(\mu_*G-\varphi(v)\,\rho^*\mathcal X_s\bigr)\big|_{E'_J},
\]

using once more that the coefficient of \(E'_j\) in \(\mu_*G\) is
\(b'_j\varphi(e'_j)\) and that
\(\rho^*\mathcal X_s=\mathcal X'_s=\sum_{j}b'_jE'_j\).  The total special
fibre is the divisor of \(\pi\), hence principal, so its restriction to
\(E'_J\) is trivial as a class.  This proves
\eqref{eq:ch22-slope-divisor-identity}.
\end{proof}

\section{The slope estimate}
\label{sec:ch22-slope}

We can now bound the derivative of a model psh function from an arbitrary
face toward a vertex.  The statement is an induction on the dimension of the
face; the base case is the vertex estimate of
\Cref{prop:ch22-vertex-estimate}, and each inductive step is one application
of the star-shrinking construction followed by one intersection inequality
on the new stratum.

\begin{proposition}
\label{prop:ch22-slope-estimate}
There is a constant \(C_{\mathrm{slope}}=C(\mathcal X,\theta,\mathcal A)\)
with the following property.  Let \(\varphi\) be a \(\theta\)-psh model
function normalized by \eqref{eq:ch22-normalization}, let \(\tau\) be a face
of \(\Delta_{\mathcal X}\), let \(e\) be a vertex of \(\tau\), let \(\sigma\)
be the face of \(\tau\) opposite \(e\), and let \(v\in\RelInt(\sigma)\).
Then

\[
                    D_v\varphi(e)\geq-C_{\mathrm{slope}} .
\tag{22.43}\label{eq:ch22-slope-lower-bound}
\]

Consequently \(\operatorname{Lip}_{\mathrm f}(\varphi|_{\Delta_{\mathcal X}})
\leq C'\) for a constant \(C'\) depending only on
\((\mathcal X,\theta,\mathcal A)\) and the chosen compatible norms.
\end{proposition}

\begin{proof}
Write \(q=\dim\tau\), so that \(\dim\sigma=q-1\) and the vertex set
\(J\) of \(\sigma\) has \(\#J=q\).  We prove
\eqref{eq:ch22-slope-lower-bound} by induction on \(q\), simultaneously with
the statement

\[
 \mathrm{(S}_q\mathrm{)}\qquad
 \begin{aligned}
 &\text{for every face }\tau'\text{ of dimension}\leq q:\\
 &\quad\operatorname{Lip}(\varphi|_{\tau'})\leq C_q
 \quad\text{and}\quad
 |\varphi|\leq C_q\ \text{on}\ \tau' ,
 \end{aligned}
\tag{22.44}\label{eq:ch22-induction-statement}
\]

with constants \(C_q\) depending only on the fixed data.

\paragraph{Base of the induction.}
For \(q=0\) the face \(\tau'\) is a vertex, the Lipschitz assertion is
empty, and the bound on \(|\varphi|\) is
\Cref{prop:ch22-vertex-estimate}; thus \(\mathrm{(S}_0\mathrm{)}\) holds
with \(C_0=C_{\mathrm{vert}}\).  The inductive step below is carried out for
every \(q\geq1\); for \(q=1\) the face \(\sigma\) is a single vertex, the
point \(v\) is that vertex, and the terms indexed by \(j\in J\) in
\eqref{eq:ch22-slope-intersection} vanish because \(e_j-v=0\).

\paragraph{The inductive step: reduction to a good \(v\).}
Assume \(\mathrm{(S}_{q-1}\mathrm{)}\), and let \(\tau,e,\sigma,v\) be as in
the statement with \(\dim\tau=q\).  The function \(\varphi\) is piecewise
affine on \(\Delta_{\mathcal X}\): it is determined by a vertical divisor on
some model dominating \(\mathcal X\), hence affine on the faces of a
refinement of \(\Delta_{\mathcal X}\), by
\Cref{prop:skeleton-model-pa}.  Call \(v\in\RelInt(\sigma)\) \emph{good} if
it is rational and lies in the relative interior of a maximal cell of the
induced subdivision of \(\sigma\).  Good points are dense in
\(\RelInt(\sigma)\).  If \(v\) is good, then for all sufficiently small
rational \(\varepsilon>0\) the function \(\varphi\) is affine on the simplex
\(\sigma'_J\) of \eqref{eq:ch22-star-shrink} and on each segment
\([v,e'_j]\), \(j\in L\): the first because a whole neighborhood of \(v\) in
\(\sigma\) lies in one cell of affineness, the second because for small
\(\varepsilon\) the segment \([v,e'_j]\) lies in the closed cell of the
refinement that the direction \(e_j-v\) enters at \(v\).

Suppose the bound \eqref{eq:ch22-slope-lower-bound} proved for all good
\(v\).  For fixed \(e\), the function
\(v\mapsto D_v\varphi(e)\) is, by
\eqref{eq:ch22-directional-derivative} and monotonicity of difference
quotients, the infimum over \(t\in(0,1]\) of the continuous functions
\(v\mapsto t^{-1}(\varphi((1-t)v+te)-\varphi(v))\), hence upper
semicontinuous on \(\RelInt(\sigma)\) by
\Cref{lem:usc-inf}.  If \(v_k\to v\) with \(v_k\) good, upper
semicontinuity gives
\(D_v\varphi(e)\geq\varlimsup_kD_{v_k}\varphi(e)\geq-C_{\mathrm{slope}}\).
So it suffices to treat good \(v\).

\paragraph{The inductive step: the intersection inequality.}
Let \(v\) be good and let \(\varepsilon\) be small enough for the affineness
statements above.  Apply \Cref{lem:ch22-star-shrinking} to obtain
\(\Delta'\) and \(\rho:\mathcal X'\to\mathcal X\), and let
\(\mu:\mathcal Y\to\mathcal X'\) dominate a determination \(G\) of
\(\varphi\), with \(\pi=\rho\circ\mu\) and \(\pi^*\theta_{\mathcal X}+[G]\)
relatively nef.  Let \(E'_J\) be the stratum of \(\mathcal X'\) attached to
\(\sigma'_J\).  Since \(\sigma'_J\subset\sigma\) has the same dimension and
\(v\in\RelInt(\sigma'_J)\), the morphism \(E'_J\to E_J\) is dominant between
varieties of the same dimension \(n+1-q\), hence generically finite, of some
degree \(\delta>0\).

The divisor \(G\) is \(\mu\)-nef: \(\pi^*\theta_{\mathcal X}+[G]\) is
relatively nef, and \(\pi^*\theta_{\mathcal X}\) has degree zero on every
curve contracted by \(\mu\), because such a curve is contracted by
\(\pi\) as well.  Since \(\mathcal X'\) is vertically
\(\mathbf Q\)-factorial, \Cref{lem:ch22-negativity} applies and

\[
 F:=\mu^*\mu_*G-G\geq0 .
\tag{22.45}\label{eq:ch22-negativity-application}
\]

Let \(W\subset\mathcal Y\) be the closure of the center of \(v\) on
\(\mathcal Y\); it maps onto \(E'_J\), the center of \(v\) on
\(\mathcal X'\), by \Cref{lem:skeleton-monomial-invariants}.  The model
functions of \(G\) and of \(\mu^*\mu_*G\) are both affine on \(\sigma'_J\)
and take the same values at its vertices \(e'_j\), \(j\in J\); hence they
agree at \(v\), so the model function of \(F\) vanishes at \(v\).  As \(F\)
is effective, \eqref{eq:ch20-sign} shows that the center of \(v\) is not
contained in \(\Supp F\); therefore \(W\not\subset\Supp F\) and \(F|_W\) is
an effective class.  Restricting the identity
\(\mu^*(\rho^*\theta_{\mathcal X}+\mu_*G)
=\pi^*\theta_{\mathcal X}+G+F\) to \(W\) gives

\[
 (\mu|_W)^*\!\left(
   (\rho^*\theta_{\mathcal X}+\mu_*G)|_{E'_J}\right)
   =(\pi^*\theta_{\mathcal X}+G)|_W+F|_W,
\tag{22.46}\label{eq:ch22-nef-effective-decomposition}
\]

a sum of a nef class and an effective class.  Now apply
\Cref{lem:ch22-pullback-positivity} with \(g=\mu|_W:W\to V:=E'_J\),
\(\alpha=(\rho^*\theta_{\mathcal X}+\mu_*G)|_{E'_J}\) and
\(\beta=\rho^*\mathcal A|_{E'_J}\), which is nef.  Since
\(\dim V=n+1-q\), the exponent in
\eqref{eq:ch22-pullback-pseudoeffective} is \(\dim V-1=n-q=n-\#J\), and we
obtain

\[
 \Bigl(\bigl(\rho^*\theta_{\mathcal X}+\mu_*G\bigr)\cdot E'_J\cdot
       (\rho^*\mathcal A)^{n-\#J}\Bigr)\geq0 .
\]

Replace \(\mu_*G|_{E'_J}\) by \eqref{eq:ch22-slope-divisor-identity} and
push forward by \(\rho\), using the projection formula
\eqref{eq:appF-projection} and \(\rho_*E'_J=\delta E_J\).  Dividing by
\(\delta>0\) yields

\[
 \sum_{j\in L}D_v\varphi(e_j)\,b_j
       \bigl(E_J\cdot E_j\cdot\mathcal A^{n-\#J}\bigr)
 \geq
 -\bigl(E_J\cdot\theta_{\mathcal X}\cdot
                         \mathcal A^{n-\#J}\bigr).
\tag{22.47}\label{eq:ch22-slope-intersection}
\]

\paragraph{Isolating the term of \(e\).}
The exponent in \eqref{eq:ch22-slope-intersection} deserves a remark, since
it corrects the printed display \cite[(6.3)]{BFJ16}: the stratum \(E_J\) has
dimension \(n+1-\#J\), so intersecting it with one divisor class and
\(\mathcal A^{k}\) gives a number exactly when \(k=n-\#J\); the printed
exponent \(n-\#J-1\) leaves a one-cycle and is negative for
\(n=\#J=1\).  With the corrected exponent, all terms of
\eqref{eq:ch22-slope-intersection} are honest intersection numbers on
\(\mathcal X\), and their absolute values are bounded by a constant
depending only on \((\mathcal X,\theta,\mathcal A)\).

We bound the individual derivatives occurring in
\eqref{eq:ch22-slope-intersection}.  For \(j\in J\), the point \(v\) and the
vertex \(e_j\) both lie in the face \(\sigma\) of dimension \(q-1\), so
\(D_v\varphi(e_j)\) is a directional derivative of
\(\varphi|_\sigma\) and \(\mathrm{(S}_{q-1}\mathrm{)}\) gives

\[
 |D_v\varphi(e_j)|\leq C_{q-1}\,|e_j-v|\leq C_{q-1}\operatorname{diam}
       (\Delta_{\mathcal X}),
 \qquad j\in J .
\tag{22.48}\label{eq:ch22-inductive-derivative-bound}
\]

For \(j\in L\setminus J\), convexity of \(\varphi\) along \([v,e_j]\) and
the normalization \(\varphi\leq0\) give the one-sided bound

\[
             D_v\varphi(e_j)
             \leq\varphi(e_j)-\varphi(v)\leq-\varphi(v)\leq C_{q-1},
\tag{22.49}\label{eq:ch22-slope-upper-bound}
\]

using \(\mathrm{(S}_{q-1}\mathrm{)}\) for the bound on \(|\varphi(v)|\).
Moreover, for \(j\in L\setminus J\) the divisor \(E_j|_{E_J}\) is a nonzero
effective divisor on \(E_J\), so

\[
 b_j\bigl(E_J\cdot E_j\cdot\mathcal A^{n-\#J}\bigr)\geq1 .
\tag{22.50}\label{eq:ch22-positive-slope-coefficient}
\]

Now isolate in \eqref{eq:ch22-slope-intersection} the term with
\(e_j=e\), which occurs because \(e\in L\setminus J\).  Moving all other
terms to the right-hand side, bounding those with \(j\in J\) by
\eqref{eq:ch22-inductive-derivative-bound} and those with
\(j\in L\setminus J\), \(j\neq e\), from above by
\eqref{eq:ch22-slope-upper-bound} multiplied by the nonnegative coefficients
\eqref{eq:ch22-positive-slope-coefficient}, and using
\eqref{eq:ch22-positive-slope-coefficient} for the isolated coefficient,
gives \(D_v\varphi(e)\geq-C\) with \(C\) depending only on
\((\mathcal X,\theta,\mathcal A)\) and \(C_{q-1}\).  Set
\(C_{\mathrm{slope}}\) to be the largest such constant over the finitely
many faces \(\tau\) of \(\Delta_{\mathcal X}\).

\paragraph{Closing the induction.}
The bound just proved verifies the hypothesis
\eqref{eq:ch22-simplex-hypothesis} of \Cref{lem:ch22-simplex} on each face
\(\tau\) of dimension \(q\), with \(\lambda=C_{\mathrm{slope}}\).  That
lemma gives
\(\operatorname{Lip}(\varphi|_\tau)\leq(q+2)\Lambda(\tau)C_{\mathrm{slope}}\)
and \(\max_\tau\varphi-\min_\tau\varphi\leq(q+1)C_{\mathrm{slope}}\);
combined with the vertex bound \eqref{eq:ch22-vertex-lower-bound} at any
vertex of \(\tau\), this proves \(\mathrm{(S}_q\mathrm{)}\) with

\[
 C_q:=\max\bigl\{(q+2)\Lambda_{\max}C_{\mathrm{slope}},\
       C_{\mathrm{vert}}+(q+1)C_{\mathrm{slope}}\bigr\},
 \qquad
 \Lambda_{\max}:=\max_\tau\Lambda(\tau).
\tag{22.51}\label{eq:ch22-inductive-constant}
\]

The induction terminates at \(q=n\), the maximal dimension of a face of
\(\Delta_{\mathcal X}\), and the final assertion of the proposition holds
with \(C'=C_n\).
\end{proof}

\begin{remark}[Reading the two estimates together]
\label{rem:ch22-estimates-together}
The vertex estimate and the slope estimate are not two halves of one
computation; they are the same computation performed on two different
models.  In both cases one intersects the relatively nef class
\(\pi^*\theta+G\) with an effective vertical cycle and with powers of a
fixed relatively ample class, and reads the result as a linear inequality
between the values of \(\varphi\) at the vertices of a complex.  What
changes is which complex: for the vertex estimate it is
\(\Delta_{\mathcal X}\) itself, and for the slope estimate it is the shrunk
star of \Cref{lem:ch22-star-shrinking}, whose vertex values encode
derivatives at \(v\).  Everything else --- the sign of the off-diagonal
terms, the vanishing of the row sums, the propagation --- is common.
\end{remark}

\section{Equicontinuity on dual complexes}
\label{sec:ch22-equicontinuity}

\begin{theorem}\label{thm:bfj-skeletal-equicontinuity}
Let \(\mathcal X\) be an SNC model on which \(\theta\) is determined.  There
is a constant \(C=C(\mathcal X,\theta)>0\) with the following properties.

\begin{enumerate}
\item For every \(\theta\)-psh function \(\varphi\), the restriction
\(\varphi|_{\Delta_{\mathcal X}}\) has Lipschitz constant at most \(C\) on
each face.
\item If \eqref{eq:ch22-normalization} holds, then

\[
                  -C\leq\varphi\leq0
                  \quad\text{on }\Delta_{\mathcal X}.
\tag{22.52}\label{eq:ch22-skeletal-bound}
\]
\end{enumerate}

In particular, normalized restrictions form a uniformly bounded,
equicontinuous family in \(C^0(\Delta_{\mathcal X})\).
\end{theorem}

\begin{proof}
Fix a relatively ample line bundle \(\mathcal A\) on \(\mathcal X\) and let
\(C\) be the maximum of the constants \(C'\) of
\Cref{prop:ch22-slope-estimate} and
\(C_{\mathrm{vert}}+C'\delta(\Delta_{\mathcal X})\), where
\(\delta(\Delta_{\mathcal X})\) is the diameter of
\Cref{lem:ch22-path-metric}.

\emph{Model psh functions.}  For a \(\theta\)-psh model function normalized
by \eqref{eq:ch22-normalization}, assertion (1) is the last statement of
\Cref{prop:ch22-slope-estimate} and assertion (2) follows from
\Cref{prop:ch22-vertex-estimate} together with
\Cref{lem:ch22-convexity-consequences}\,\cref{itm:ch22-cvx:osc}: the maximum
is \(0\) by \eqref{eq:ch22-vertex-max}, and the facewise Lipschitz bound
propagates it along the path metric.  For a model psh function that is not
normalized, subtract its supremum; assertion (1) is unaffected by adding a
constant.

\emph{General psh functions.}  Let \(\varphi\in\PSH(X,\theta)\).  By clause
(2) of \Cref{def:ch21-bfj-psh}, the restriction
\(\varphi|_{\Delta_{\mathcal X}}\) is a uniform limit of restrictions
\(\psi_k|_{\Delta_{\mathcal X}}\) of \(\theta\)-psh model functions
\(\psi_k\).  Each \(\psi_k-\sup\psi_k\) satisfies the model case, so
\(\operatorname{Lip}_{\mathrm f}(\psi_k|_{\Delta_{\mathcal X}})\leq C\); a
uniform limit of \(C\)-Lipschitz functions is \(C\)-Lipschitz, which is (1).

For (2), assume \(\sup_{X^{\An}}\varphi=0\).  By
\Cref{prop:ch21-sup-on-skeleton} the supremum of each \(\psi_k\) is its
maximum on \(\Delta_{\mathcal X}\); hence

\[
 \bigl|\sup\nolimits_{X^{\An}}\psi_k\bigr|
 =\bigl|\max\nolimits_{\Delta_{\mathcal X}}\psi_k
    -\max\nolimits_{\Delta_{\mathcal X}}\varphi\bigr|
 \leq\|\psi_k-\varphi\|_{\Delta_{\mathcal X}}\longrightarrow0 .
\tag{22.53}\label{eq:ch22-normalizing-constants-vanish}
\]

Applying the normalized model bound to \(\psi_k-\sup\psi_k\) and letting
\(k\to\infty\) gives \(-C\leq\varphi\leq0\) on
\(\Delta_{\mathcal X}\).  The final sentence of the theorem is the
combination of (1), (2) and \Cref{lem:ch22-path-metric}.
\end{proof}

The proof produced more than the statement, and the surplus is used often
enough to be recorded separately.

\begin{corollary}
\label{cor:ch22-explicit-modulus}
With \(C\) as in \Cref{thm:bfj-skeletal-equicontinuity}, every
\(\varphi\in\PSH(X,\theta)\) satisfies

\[
 \bigl|\varphi(x)-\varphi(y)\bigr|\leq C\,d_{\Delta_{\mathcal X}}(x,y),
 \qquad x,y\in\Delta_{\mathcal X},
\tag{22.54}\label{eq:ch22-modulus}
\]

and, if \eqref{eq:ch22-normalization} holds,
\(\operatorname{osc}_{\Delta_{\mathcal X}}\varphi
\leq C\,\delta(\Delta_{\mathcal X})\).  The modulus in
\eqref{eq:ch22-modulus} is a genuine modulus of continuity for the compact
metric space \((\Delta_{\mathcal X},d_{\Delta_{\mathcal X}})\), independent
of \(\varphi\).
\end{corollary}

\begin{proof}
Combine \Cref{thm:bfj-skeletal-equicontinuity}\,(1) with
\Cref{lem:ch22-path-metric}.
\end{proof}

\begin{corollary}
\label{cor:ch22-arzela-ascoli}
For every SNC model \(\mathcal X\) determining \(\theta\), the set

\[
 \bigl\{\varphi|_{\Delta_{\mathcal X}}:
        \varphi\in\PSH(X,\theta),\ \sup\nolimits_{X^{\An}}\varphi=0\bigr\}
\]

is relatively compact in \(C^0(\Delta_{\mathcal X})\) for the uniform norm,
and its closure is a compact metrizable space.
\end{corollary}

\begin{proof}
The family is uniformly bounded and equicontinuous by
\Cref{thm:bfj-skeletal-equicontinuity} and
\Cref{cor:ch22-explicit-modulus}.  The dual complex is a compact metric
space by \Cref{lem:ch22-path-metric}, so the Arzel\`a--Ascoli theorem
applies; \(C^0\) of a compact metric space is a separable Banach space, so
the closure is compact metrizable.
\end{proof}

\begin{warning}{One complex at a time}
\Cref{cor:ch22-arzela-ascoli} is a statement about \emph{one} dual complex,
and on one dual complex sequences are legitimate: the space is compact
metrizable, so every sequence of normalized restrictions has a uniformly
convergent subsequence.  What is not legitimate is to perform such an
extraction simultaneously for all models by a diagonal argument, because the
model system is not countably cofinal.  This is the exact division of labour
underlying the fixed-skeleton selection principle
\Cref{prop:ch23-fixed-skeleton-selection} and its warning, and it is
quantified in \Cref{ex:ch22-countable-blind-spot}.
\end{warning}

\begin{proposition}
\label{prop:ch22-characteristic-free-estimates}
Let \(K\) be a complete discretely valued field with \(v(\pi)=1\), of
arbitrary residue characteristic, let \(X\) be a smooth connected projective
\(K\)-variety of dimension \(n\), let \(\theta\) be a closed model
\((1,1)\)-form, and let \(\mathcal X\) be an SNC model on which
\(\theta\) is determined.  Then
\Cref{prop:ch22-vertex-estimate,prop:ch22-slope-estimate} hold verbatim for
all \(\theta\)-psh model functions, hence so do assertions (1) and (2) of
\Cref{thm:bfj-skeletal-equicontinuity} for every function whose restriction
to \(\Delta_{\mathcal X}\) is a uniform limit of restrictions of
\(\theta\)-psh model functions.
\end{proposition}

\begin{proof}
We list the inputs of the two propositions and note that each is available
without a hypothesis on the residue characteristic.  The signs and the row
sum in \Cref{lem:ch22-incidence-signs} are intersection theory on the
projective \(R\)-scheme \(\mathcal X\); connectedness of \(\mathcal X_s\)
follows from the theorem on formal functions, as in the proof of
\Cref{thm:ch20-ddc-exact}.  The projection formula and Kleiman's criterion
are \Cref{thm:appF-projection} and \cite[Section~1.4]{KM08}.  The negativity
lemma \Cref{lem:ch22-negativity} was proved above using only relative global
generation.  The realization of a rational projective simplicial subdivision
by a normal, vertically \(\mathbf Q\)-factorial vertical blowup is the
normalized-blowup construction in the proof of
\Cref{thm:skeleton-subdivision}, which is toroidal and characteristic free;
the further resolution to an SNC model, which is where that theorem invokes
residue characteristic zero, is not used in
\Cref{lem:ch22-star-shrinking} or in
\Cref{prop:ch22-slope-estimate}.  Finally,
\Cref{lem:ch22-simplex} and \Cref{lem:ch22-path-metric} are convex geometry.
\end{proof}

\begin{remark}[What \Cref{prop:ch22-characteristic-free-estimates} does not say]
\label{rem:ch22-characteristic-free-caveat}
It does not say that the BFJ theory works in positive residue
characteristic.  It says that the \emph{estimate} on a given complex does.
Two things are still missing outside residue characteristic zero: SNC models
need not be cofinal, so \(\Delta_{\mathcal X}\) need not see enough of
\(X^{\An}\) for \Cref{thm:bfj-skeletal-characterization} to reconstruct
\(\varphi\); and the regularization theorem of
\Cref{sec:ch22-regularization} is unavailable, since its proof runs through
multiplier ideals.  The precise inventory is
\Cref{sec:ch22-char-zero}.
\end{remark}

\section{Two moduli computed}
\label{sec:ch22-worked-moduli}

The constants of \Cref{prop:ch22-vertex-estimate,prop:ch22-slope-estimate}
are explicit but generous.  This section computes both the general constants
and the sharp ones in the two smallest cases: a segment, where the sharp
constant can be found by convexity alone, and a triangle, where the gauge
\eqref{eq:ch22-gauge-triangle} is attained.

\begin{example}[The two-vertex estimate]
\label{ex:ch22-two-vertex}
Let \(X\) be a smooth projective curve whose reduced special fibre on a
regular semistable model \(\mathcal X\) has two components meeting
transversally at one residue-field-rational point, so that
\(b_0=b_1=1\), \((E_0\cdot E_1)=1\) and
\((E_i^2)=-1\) by \eqref{eq:ch20-two-component-matrix}.  Write

\[
              d_i=(\theta\cdot E_i),\qquad i=0,1,
\]

so that \(d_0+d_1=\deg_XL\) by \eqref{eq:ch20-special-fibre-degree}.  For
vertex values \(u_0,u_1\), formula
\eqref{eq:ch20-curve-nef-test} becomes

\[
                   -d_0\leq u_1-u_0\leq d_1.
\tag{22.55}\label{eq:ch22-two-vertex-bound}
\]

If \(\max\{u_0,u_1\}=0\), then

\[
                   -\max\{d_0,d_1\}\leq u_i\leq0.
\]

This is \eqref{eq:ch22-skeletal-bound} in its smallest nontrivial form, and
it agrees exactly with the general constant: here \(n=1\), so
\(\mathcal A\) does not appear, \(w_{01}=w_{10}=1\), \(\lambda=1\),
\(c=\max\{d_0,d_1\}\), \(\#I=2\), and
\eqref{eq:ch22-vertex-lower-bound} gives
\(C_{\mathrm{vert}}=c\).  The higher-dimensional vertex proof is the same
graph propagation with intersection weights.
\end{example}

\begin{example}[A modulus on a segment, with numbers]
\label{ex:ch22-modulus-segment}
Keep the model of \Cref{ex:ch22-two-vertex} and take

\[
              d_0=1,\qquad d_1=2,\qquad \deg_XL=3 .
\]

Parametrize \(\Delta_{\mathcal X}\) by the lattice-length coordinate
\(t\in[0,1]\) of \Cref{ex:ch22-edge-normalization}, with \(t=0\) at
\(e_0\).

\emph{The nef test in slope form.}  Let \(u=\varphi|_{\Delta_{\mathcal X}}\)
for a \(\theta\)-psh model function \(\varphi\) determined on a refinement
subdividing \([0,1]\) at rational points.  In such a refinement the
multiplicities of the components and the lattice length of an adjacent edge
are matched so that, for two adjacent components \(F,F'\) at parameters
\(t_F<t_{F'}\),

\[
 b_Fb_{F'}\bigl(F\cdot F'\bigr)=\frac1{t_{F'}-t_F} .
\]

For the single blowup of the node this is the computation of
\Cref{ex:refinement-segment}: there \(b_F=2\),
\(b_{E'_0}=b_{E'_1}=1\), each adjacent intersection number is one, and
each new edge has lattice length \(\tfrac12\).  Thus the product of the two
multiplicities is \(2=1/\tfrac12\) on both edges.  The general case follows
by iterating, as in \Cref{rem:lattice-length}.

Consequently, after multiplying the inequality at \(E_i\) in
\eqref{eq:ch20-curve-nef-test} by the positive number \(b_i\), the term
\[
 b_i\sum_{j\neq i}b_j(E_j\cdot E_i)(u_j-u_i)
\]
is exactly the sum of the outgoing slopes of \(u\) at that vertex, computed
in the coordinate \(t\); the reference term becomes \(b_id_i\).  On the
present refinement the two original components have multiplicity one, while
every inserted component has reference degree zero by the projection
formula.  Hence these positive rescalings do not change the following nef
test, which reads

\[
 \begin{aligned}
 \text{at an interior vertex:}&\quad
   u'(t_k^+)-u'(t_k^-)\geq0,\\
 \text{at }e_0:&\quad d_0+u'(0^+)\geq0,\\
 \text{at }e_1:&\quad d_1-u'(1^-)\geq0 .
 \end{aligned}
\tag{22.56}\label{eq:ch22-curve-slope-form}
\]

The first line is convexity; the other two are the boundary conditions.
Hence

\[
 -1=-d_0\leq u'(0^+)\leq u'(t)\leq u'(1^-)\leq d_1=2
\]

for every \(t\), so \(\operatorname{Lip}(u)\leq2\) and, by
\eqref{eq:ch22-vertex-lower-bound}, \(-2\leq u\leq0\) after normalization.
In the language of \Cref{lem:ch22-simplex}, hypothesis
\eqref{eq:ch22-simplex-hypothesis} holds with
\(\lambda=\max\{d_0,d_1\}=2\), the gauge is \(\Lambda=1\) by
\eqref{eq:ch22-gauge-segment}, and the sharp one-dimensional form of
\Cref{rem:ch22-simplex-constant} gives \(\operatorname{Lip}(u)\leq2\), which
is what the boundary conditions gave directly.

\emph{The extremal function.}  Put

\[
 u_\star(t):=\max\{-d_0t,\ d_1(t-1)\}=\max\{-t,\ 2t-2\}.
\tag{22.57}\label{eq:ch22-extremal-segment}
\]

Then \(u_\star(0)=u_\star(1)=0\), the two slopes are \(-1\) and \(2\), the
two branches cross at

\[
 t_\star=\frac{d_1}{d_0+d_1}=\frac23,
 \qquad
 \min u_\star=u_\star(t_\star)=-\frac{d_0d_1}{d_0+d_1}=-\frac23 ,
\]

and the slope jumps upward by \(d_0+d_1=3\) there, so \(u_\star\) is convex
and both boundary conditions in \eqref{eq:ch22-curve-slope-form} hold with
equality.  By \Cref{prop:ch20-model-max} and
\eqref{eq:ch20-nodal-ideal}, \(u_\star\) is the restriction of the model
function \(\log|(x,y^2)|\) of a vertical ideal on the node chart
\(xy=\pi\), and it is \(\theta\)-psh.

\emph{Both general constants are attained.}  Consider instead the affine
function

\[
 u_\flat(t):=d_1(t-1)=2t-2 .
\tag{22.58}\label{eq:ch22-flat-extremal}
\]

It is convex, its maximum is \(0\) at \(t=1\), and the two boundary
conditions of \eqref{eq:ch22-curve-slope-form} read \(1+2\geq0\) and
\(2-2=0\), the second with equality.  Hence \(u_\flat\) is the restriction
of a normalized \(\theta\)-psh model function, and

\[
 \min u_\flat=-d_1=-2=-C_{\mathrm{vert}},
 \qquad
 \operatorname{Lip}(u_\flat)=2=\max\{d_0,d_1\}.
\tag{22.59}\label{eq:ch22-segment-sharp-constants}
\]

So in this example both the vertex bound
\eqref{eq:ch22-vertex-lower-bound} and the facewise Lipschitz bound are
equalities: the general constants cannot be improved without further
hypotheses.  The function \(u_\star\) is extremal in a complementary sense:
among normalized restrictions vanishing at \emph{both} vertices, it is the
deepest, with

\[
 \min u_\star=-\frac{d_0d_1}{d_0+d_1}=-\frac23 ,
\]

both boundary conditions holding with equality.  The contrast is worth
remembering.  The vertex estimate is governed by the worst single vertex and
gives \(\max\{d_0,d_1\}\); the two-sided extremal is governed by a
harmonic-type mean and gives \(d_0d_1/(d_0+d_1)\), which is three times
smaller here.  A normalization that pins the function at both ends is
therefore strictly stronger than the normalization
\eqref{eq:ch22-normalization}.
\end{example}

\begin{example}[A modulus on a triangle, with numbers]
\label{ex:ch22-modulus-triangle}
Let \(n=2\) and let \(\mathcal X\) contain the chart
\(\mathcal U=\Spec R[x,y,z]/(xyz-\pi)\) of
\Cref{ex:toric-triangle-chart}, with three components of multiplicity one
meeting pairwise in curves and all three in one point, so that the
corresponding face of \(\Delta_{\mathcal X}\) is the standard triangle
\(\sigma\) of \Cref{ex:dual-complex-triangle}.  Use the coordinates
\((t_1,t_2)=(\alpha_x,\alpha_y)\), in which \(\sigma\) is
\(\operatorname{conv}\{(0,0),(1,0),(0,1)\}\), and the compatible Euclidean
norm of \Cref{ex:ch22-gauge-computed}, so that
\(\Lambda(\sigma)=\sqrt2\).

\emph{An explicit convex model function.}  Let

\[
\mathfrak a:=\bigl(xy,\ yz,\ zx\bigr)^3
\]

and \(u:=\log|\mathfrak a|\).  By \eqref{eq:ch20-max-ideal-function} and
\Cref{lem:vertical-ideal-convex},

\[
 u=3\log|(xy,yz,zx)|=\max\{\ell_0,\ell_1,\ell_2\},
 \qquad
 \begin{aligned}
 \ell_0&=-3(\alpha_x+\alpha_y)=-3t_1-3t_2,\\
 \ell_1&=-3(\alpha_y+\alpha_z)=-3+3t_1,\\
 \ell_2&=-3(\alpha_x+\alpha_z)=-3+3t_2 .
 \end{aligned}
\tag{22.60}\label{eq:ch22-triangle-function}
\]

The three affine functions sum to \(-6=-3\,v(\pi^2)\), consistently with
\((xy)(yz)(zx)=\pi^2\).  By
\Cref{lem:ch20-generated-ideal-psh}, \(u\) is \(\theta\)-psh for any
\(\theta\) that is the curvature of a model metric on a line bundle
\(\mathcal L\) with \(\mathcal L\otimes\mathfrak a\) relatively globally
generated.

\emph{The numbers.}  At the three vertices \(u=0\); at the barycenter all
three branches equal \(-2\); the gradients are
\(\nabla\ell_0=(-3,-3)\) and \(\nabla\ell_1=(3,0)\),
\(\nabla\ell_2=(0,3)\).  Hence

\[
 \operatorname{osc}_\sigma u=2,
 \qquad
 \operatorname{Lip}(u)=|\nabla\ell_0|=3\sqrt2\approx4.243 .
\tag{22.61}\label{eq:ch22-triangle-numbers}
\]

\emph{The hypothesis constant.}  We compute the smallest \(\lambda\) for
which \eqref{eq:ch22-simplex-hypothesis} holds.  For \(v=(s,1-s)\) on the
edge opposite \(e_0=(0,0)\) one has, for \(0<s<1\),

\[
 D_vu(e_0)=
 \begin{cases}
  -3(1-s)&\text{if }s<\tfrac12,\\
  -3s&\text{if }s>\tfrac12,
 \end{cases}
\]

with the value \(-\tfrac32\) at \(s=\tfrac12\); the infimum over the open
edge is \(-3\), approached as \(v\) tends to a vertex.  The same computation
at the other two vertices gives the same value.  Hence
\(\lambda=3\), and

\[
 \operatorname{Lip}(u)=3\sqrt2=\Lambda(\sigma)\lambda ,
\tag{22.62}\label{eq:ch22-triangle-sharp}
\]

so the gauge factor of \Cref{lem:ch22-simplex} is attained, exactly as
\Cref{rem:ch22-simplex-constant} predicts.  The general bound of that lemma
reads \((q+2)\Lambda(\sigma)\lambda=4\cdot\sqrt2\cdot3=12\sqrt2\approx16.97\),
and its oscillation bound reads \((q+1)\lambda=9\); the true values are
\(3\sqrt2\) and \(2\).  The dimensional factor is therefore lossy, but the
directional factor is not.

\emph{Where the numbers come from geometrically.}  The value \(\lambda=3\)
is the exponent chosen in \(\mathfrak a\); replacing
\(\mathfrak a\) by \((xy,yz,zx)^k\) multiplies both \(\lambda\) and
\(\operatorname{Lip}(u)\) by \(k/3\), so the ratio
\eqref{eq:ch22-triangle-sharp} is a property of the triangle and not of the
ideal.  The constant \(\Lambda(\sigma)=\sqrt2\) is the price of measuring
lengths with the Euclidean norm of
\Cref{def:ch22-compatible-norm}; with the \(\ell^1\)-norm of the barycentric
coordinates the same computation gives the constant \(1\).  This is the
precise sense in which the choice of compatible norms changes constants and
nothing else.
\end{example}

\section{The BFJ topology and compactness}
\label{sec:ch22-topology}

\subsection{The topology}

\begin{definition}\label{def:ch22-bfj-topology}
The \emph{BFJ weak topology} on \(\PSH(X,\theta)\) is the topology of uniform
convergence on every dual complex.  Equivalently, it is the subspace
topology for the map

\[
 \PSH(X,\theta)\longrightarrow
 \prod_{\mathcal X}C^0(\Delta_{\mathcal X}),
 \qquad
 \varphi\longmapsto
       (\varphi|_{\Delta_{\mathcal X}})_{\mathcal X},
\tag{22.63}\label{eq:ch22-product-embedding}
\]

the product being over all SNC models on which \(\theta\) is determined,
each factor carrying the uniform norm.
\end{definition}

Three comments fix the meaning of the definition.

First, the index set of \eqref{eq:ch22-product-embedding} is directed by
domination, and by \Cref{prop:snc-cofinal} the SNC models on which
\(\theta\) is determined are cofinal among all models; so nothing is lost by
restricting to them, and by \Cref{prop:ch21-cofinal-independence} nothing
would be gained by enlarging the family.  This is the first, and the most
easily overlooked, place where residue characteristic zero is used in this
chapter: not in an estimate, but in the assertion that the coordinates of
\eqref{eq:ch22-product-embedding} see enough of \(X^{\An}\).

Second, the map \eqref{eq:ch22-product-embedding} is injective.  Indeed
\Cref{thm:bfj-skeletal-characterization}\,(2) reconstructs \(\varphi\) from
its skeletal restrictions through
\eqref{eq:ch21-skeletal-reconstruction}.  Hence the BFJ topology is
Hausdorff, being the initial topology of an injective map into a Hausdorff
space.

Third, the product may be uncountable, and then it is not metrizable.  Thus
compactness naturally gives convergent subnets, not necessarily
subsequences, and every statement below is phrased accordingly; see
\Cref{prop:net-criterion}\,(3) and \Cref{rem:weak-metrizability}.

\begin{lemma}
\label{lem:ch22-convergence-criterion}
A net \((\varphi_\alpha)\) in \(\PSH(X,\theta)\) converges to \(\varphi\) in
the BFJ topology if and only if
\(\varphi_\alpha|_{\Delta_{\mathcal X}}\to\varphi|_{\Delta_{\mathcal X}}\)
uniformly for every SNC model \(\mathcal X\) determining \(\theta\).  If the
net is decreasing and \(\varphi=\inf_\alpha\varphi_\alpha\) is not
identically \(-\infty\), then it converges in the BFJ topology.
\end{lemma}

\begin{proof}
The first assertion is the definition of the initial topology together with
\Cref{prop:net-criterion}\,(2).  For the second, fix a determining SNC model
\(\mathcal X\).  By \Cref{thm:ch21-decreasing-stability} the limit
\(\varphi\) is \(\theta\)-psh, so its restriction to
\(\Delta_{\mathcal X}\) is continuous; the restrictions
\(\varphi_\alpha|_{\Delta_{\mathcal X}}\) form a decreasing net of
continuous functions on a compact space with continuous pointwise limit, so
Dini's theorem for nets, \Cref{prop:dini-nets} and
\Cref{thm:appE-dini}, gives uniform convergence.
\end{proof}

\subsection{Properness and compactness}

\begin{theorem}\label{thm:ch22-proper-normalization}
For every model function \(\rho\), the map

\[
 \PSH(X,\theta)\longrightarrow\mathbf R,
 \qquad
 \varphi\longmapsto\sup_{X^{\An}}(\varphi-\rho)
\tag{22.64}\label{eq:ch22-proper-map}
\]

is continuous and proper.  The BFJ topology is also the topology of
pointwise convergence on either all divisorial points or all quasimonomial
points.
\end{theorem}

\begin{proof}
\emph{Step 1: the supremum is skeletal.}  Choose one SNC model
\(\mathcal X\) determining \(\theta\) and \(\rho\).  Since
\(\rho=\rho\circ p_{\mathcal X}\) by
\Cref{lem:ch21-model-skeletal-properties} and
\(\varphi\leq\varphi\circ p_{\mathcal X}\) by
\eqref{eq:ch21-retraction-inequality},

\[
 \sup_{X^{\An}}(\varphi-\rho)
 =\max_{\Delta_{\mathcal X}}(\varphi-\rho).
\tag{22.65}\label{eq:ch22-proper-map-skeletal}
\]

Uniform convergence on this one skeleton therefore proves continuity of
\eqref{eq:ch22-proper-map}: the maximum over a fixed compact set is a
\(1\)-Lipschitz function of the restriction for the uniform norm.

\emph{Step 2: uniform bounds on a sublevel set.}  Fix \(A>0\) and let
\(\mathcal F_A\) be the set of \(\varphi\) with
\(|\sup(\varphi-\rho)|\leq A\).  Since

\[
 \bigl|\sup(\varphi-\rho)-\sup\varphi\bigr|
       \leq\|\rho\|_{\infty},
\]

the constants \(\sup\varphi\) are bounded by
\(A+\|\rho\|_\infty\) on \(\mathcal F_A\).  Subtracting them and applying
\Cref{thm:bfj-skeletal-equicontinuity} and
\Cref{cor:ch22-explicit-modulus}, we find that for every SNC model
\(\mathcal Y\) the set of restrictions
\(\{\varphi|_{\Delta_{\mathcal Y}}:\varphi\in\mathcal F_A\}\) is uniformly
bounded and equicontinuous in \(C^0(\Delta_{\mathcal Y})\).  Let
\(K_{\mathcal Y}\) be the closure of that set; it is compact by
Arzel\`a--Ascoli, as in \Cref{cor:ch22-arzela-ascoli}.

\emph{Step 3: the product is compact.}  The product
\(\prod_{\mathcal Y}K_{\mathcal Y}\) is compact by Tychonoff's theorem
\cite[Chapter~5]{Kel75}, and the image of \(\mathcal F_A\) under
\eqref{eq:ch22-product-embedding} lies in it.  By
\Cref{def:ch22-bfj-topology}, \eqref{eq:ch22-product-embedding} is a
homeomorphism onto its image, so it suffices to prove that the image of
\(\mathcal F_A\) is closed in \(\prod_{\mathcal Y}K_{\mathcal Y}\).

\emph{Step 4: identifying a limit point.}  Let
\((f_{\mathcal Y})_{\mathcal Y}\) lie in the closure of that image, and
choose a net \((\varphi_\beta)\) in \(\mathcal F_A\) with
\(\varphi_\beta|_{\Delta_{\mathcal Y}}\to f_{\mathcal Y}\) uniformly for
every \(\mathcal Y\); such a net exists by
\Cref{prop:net-criterion}\,(1) applied in the product.  Uniform limits
preserve the two compatibility relations

\[
 \begin{aligned}
 f_{\mathcal Z}|_{\Delta_{\mathcal Y}}&=f_{\mathcal Y},\\
 f_{\mathcal Z}&\leq f_{\mathcal Y}\circ p_{\mathcal Y}
                    \quad\text{on }\Delta_{\mathcal Z},
 \end{aligned}
 \qquad \mathcal Z\geq\mathcal Y,
\tag{22.66}\label{eq:ch22-product-compatibility}
\]

the first because \(\Delta_{\mathcal Y}\subset\Delta_{\mathcal Z}\) by
\Cref{prop:skeleton-nesting}, the second because each \(\varphi_\beta\)
satisfies \eqref{eq:ch21-retraction-inequality}.  Define

\[
           f(x):=\inf_{\mathcal Y}
                  f_{\mathcal Y}(p_{\mathcal Y}(x)),
           \qquad x\in X^{\An}.
\tag{22.67}\label{eq:ch22-product-limit-function}
\]

Each \(f_{\mathcal Y}\circ p_{\mathcal Y}\) is continuous, so \(f\) is upper
semicontinuous by \Cref{lem:usc-inf}.  If
\(x\in\Delta_{\mathcal Y}\), the first relation in
\eqref{eq:ch22-product-compatibility} shows that every term of the infimum
indexed by \(\mathcal Z\geq\mathcal Y\) equals \(f_{\mathcal Y}(x)\), and
the family is directed, so \(f|_{\Delta_{\mathcal Y}}=f_{\mathcal Y}\).  The
second relation gives \(f\leq f\circ p_{\mathcal Y}\), which is
\eqref{eq:ch21-retraction-inequality}.  Each \(f_{\mathcal Y}\) belongs to
the uniform closure of the restrictions of \(\theta\)-psh model functions,
because each \(\varphi_\beta\) does and that closure is closed; this is
clause (2) of \Cref{def:ch21-bfj-psh}.  Since every skeletal restriction is
finite, \(f\not\equiv-\infty\).  Thus \(f\in\PSH(X,\theta)\), and

\[
 \sup_{X^{\An}}(f-\rho)=\max_{\Delta_{\mathcal X}}(f-\rho)
 =\lim_\beta\max_{\Delta_{\mathcal X}}(\varphi_\beta-\rho)\in[-A,A]
\]

by \eqref{eq:ch22-proper-map-skeletal}.  Hence \(f\in\mathcal F_A\), its
image is the given point, and \(\mathcal F_A\) is compact.  This proves
properness.

\emph{Step 5: the two coarser topologies.}  Uniform convergence on every
dual complex implies pointwise convergence at every quasimonomial point,
since each such point lies on some dual complex by
\Cref{thm:skeleton-embedding}, and a fortiori at every divisorial point.
Conversely, let \((\varphi_\alpha)\) converge pointwise at all divisorial
points to some \(\varphi\in\PSH(X,\theta)\), fix a determining SNC model
\(\mathcal X\), and let \(\varepsilon>0\).  Pointwise convergence at one
divisorial vertex bounds the constants \(\sup\varphi_\alpha\) beyond some
index, by \eqref{eq:ch22-skeletal-bound} applied to
\(\varphi_\alpha-\sup\varphi_\alpha\); hence
\Cref{cor:ch22-explicit-modulus} provides one modulus \(C\) valid for all
sufficiently advanced \(\alpha\) and for \(\varphi\).  Choose a finite set
\(S\subset\Delta_{\mathcal X}\) of divisorial points that is
\(\varepsilon/(3C)\)-dense for \(d_{\Delta_{\mathcal X}}\); this is possible
because divisorial points are dense in every dual complex and the complex is
compact.  Pointwise convergence on the finite set \(S\) is uniform on
\(S\) beyond one index, and the common modulus gives, for
\(x\in\Delta_{\mathcal X}\) and \(y\in S\) with
\(d(x,y)\leq\varepsilon/(3C)\),

\[
 |\varphi_\alpha(x)-\varphi(x)|
 \leq|\varphi_\alpha(x)-\varphi_\alpha(y)|
   +|\varphi_\alpha(y)-\varphi(y)|
   +|\varphi(y)-\varphi(x)|
 <\varepsilon .
\]

Thus divisorial pointwise convergence implies uniform convergence on every
dual complex.  The same argument, with the density step omitted, applies to
pointwise convergence on all quasimonomial points.
\end{proof}

\begin{corollary}\label{thm:bfj-psh-compactness}
The normalized space

\[
 \PSH_0(X,\theta)
 :=\{\varphi\in\PSH(X,\theta):
                    \sup_{X^{\An}}\varphi=0\}
\tag{22.68}\label{eq:ch22-normalized-space}
\]

is compact.  Equivalently, \(\PSH(X,\theta)/\mathbf R\) is compact.
\end{corollary}

\begin{proof}
Take \(\rho=0\) in \Cref{thm:ch22-proper-normalization}.  The zero fibre of
the continuous proper map \(\varphi\mapsto\sup\varphi\) is compact, and
every class modulo constants has a unique representative in that fibre.
\end{proof}

\begin{corollary}
\label{cor:ch22-net-compactness}
Every net in \(\PSH(X,\theta)\) whose supremum functional
\(\varphi\mapsto\sup_{X^{\An}}\varphi\) is bounded admits a subnet
converging in the BFJ topology to some \(\psi\in\PSH(X,\theta)\); the
convergence is uniform on every dual complex.  If in addition the net is
contained in \(\PSH_0(X,\theta)\), then \(\psi\in\PSH_0(X,\theta)\).
\end{corollary}

\begin{proof}
A bounded supremum places the net in a set \(\mathcal F_A\), which is
compact by \Cref{thm:ch22-proper-normalization}.  Apply
\Cref{prop:net-criterion}\,(3), or \Cref{thm:appE-net-compactness}, to
extract a convergent subnet, and \Cref{lem:ch22-convergence-criterion} to
identify the mode of convergence.  The last assertion holds because
\(\PSH_0\) is the fibre of a continuous map, hence closed.
\end{proof}

\begin{remark}[Three topologies, and one that is strictly finer]
\label{rem:ch22-topology-comparison}
On \(\PSH(X,\theta)\) we now have three descriptions of one topology:
uniform convergence on all dual complexes, pointwise convergence on
quasimonomial points, and pointwise convergence on divisorial points.  A
fourth candidate, uniform convergence on \(X^{\An}\), is strictly finer and
is \emph{not} the BFJ topology: \Cref{ex:ch22-countable-blind-spot} produces
a net converging in the BFJ topology whose members all remain at uniform
distance \(1\) from the limit.  The distinction matters in
\Cref{chap:bfj-energy}, where the energy functional is proved to be upper
semicontinuous for the BFJ topology and continuous along decreasing nets;
those are two different statements, and neither of them is continuity for
the BFJ topology.
\end{remark}

\section{Nets, sequences, and a countable blind spot}
\label{sec:ch22-nets}

Compactness has just been proved in the form: every net with bounded
supremum has a convergent subnet.  A reader trained on metric spaces will
want to replace ``net'' by ``sequence'' and ``subnet'' by ``subsequence'',
and the purpose of this section is to show, on one completely explicit
family, what that replacement would cost.  The construction is elementary:
it uses only the blowup of a closed point of the special fibre of a curve.

\subsection{Bumps over closed points of the special fibre}

\begin{lemma}
\label{lem:ch22-bump}
Let \(k\) be an algebraically closed field of characteristic zero,
\(R=k[[t]]\), \(K=k((t))\), \(X=\mathbf P^1_K\), and let
\(\mathcal X=\mathbf P^1_R\) with special fibre \(E=\mathcal X_s\cong
\mathbf P^1_k\), a reduced irreducible component.  Let \(L=\cO_X(1)\) and let
\(\theta\) be the curvature of the model metric determined by
\(\cO_{\mathcal X}(1)\), so that \((\theta\cdot E)=1\).  For a closed point
\(z\in E(k)\) let \(\rho_z:\mathcal X_z\to\mathcal X\) be the blowup of
\(z\), with exceptional divisor \(F_z\), and put

\[
 \varphi_z:=\varphi_{-F_z}\in\mathcal D(X).
\tag{22.69}\label{eq:ch22-bump-function}
\]

Then:
\begin{enumerate}
\item\label{itm:ch22-bump:psh} \(\varphi_z\) is a \(\theta\)-psh model
function with
\[
 \sup_{X^{\An}}\varphi_z=0,
 \qquad
 \min_{X^{\An}}\varphi_z=\varphi_z(x_{F_z})=-1,
 \qquad
 \|\varphi_z\|_\infty=1 ;
\]
\item\label{itm:ch22-bump:support} \(\varphi_z(x)\neq0\) only if the center
of \(x\) on \(\mathcal X\) is \(z\);
\item\label{itm:ch22-bump:invisible} if \(\mathcal Y\) is a model of \(X\)
over \(R\) no component of whose special fibre has center \(z\) on
\(\mathcal X\), then \(\varphi_z\) vanishes identically on
\(\Delta_{\mathcal Y}\);
\item\label{itm:ch22-bump:centers} for a fixed model \(\mathcal Y\), the set
of \(z\in E(k)\) that occur as centers of components of
\(\mathcal Y_s\) is finite.
\end{enumerate}
\end{lemma}

\begin{proof}
\cref{itm:ch22-bump:psh}  The special fibre of \(\mathcal X_z\) is
\(E'+F_z\), where \(E'\) is the strict transform of \(E\); both components
are reduced, \((F_z^2)=-1\), \((E'\cdot F_z)=1\), and the dual complex is
the segment with vertices \(x_{E'}\) and \(x_{F_z}\).  By
\eqref{eq:ch20-sign} the model function of \(-F_z\) takes the values
\(0\) at \(x_{E'}\) and \(-1\) at \(x_{F_z}\).  The reference degrees are
\((\rho_z^*\theta\cdot E')=(\theta\cdot E)=1\) and
\((\rho_z^*\theta\cdot F_z)=0\), so the nef test
\eqref{eq:ch20-curve-nef-test} reads

\[
 \text{at }E':\ 1+(-1-0)=0\geq0,
 \qquad
 \text{at }F_z:\ 0+(0-(-1))=1\geq0 .
\]

Hence \(\varphi_z\) is \(\theta\)-psh at model level, and by
\Cref{thm:ch21-model-consistency} it lies in \(\PSH(X,\theta)\).  Its
supremum is its maximum on the segment,
\Cref{prop:ch21-sup-on-skeleton}, which is \(0\), and its minimum is
\(-1\).

\cref{itm:ch22-bump:support}  Away from \(F_z\) the sheaf
\(\cO_{\mathcal X_z}(-F_z)\) is trivialized by the constant \(1\), so
\(\varphi_z\) vanishes at every point whose center on \(\mathcal X_z\) is
not contained in \(F_z\).  Since \(\rho_z\) is an isomorphism outside
\(F_z\) and contracts \(F_z\) to \(z\), such points are exactly those whose
center on \(\mathcal X\) differs from \(z\).

\cref{itm:ch22-bump:invisible}  Let \(x\in\Delta_{\mathcal Y}\).  By
\Cref{lem:skeleton-monomial-invariants} the center of \(x\) on
\(\mathcal Y\) is the generic point of the stratum \(E_J\) of
\(\mathcal Y_s\) whose face contains \(x\) in its relative interior, and the
center of \(x\) on \(\mathcal X\) is the image of that generic point.  That
image is contained in the image of each component \(E_j\), \(j\in J\).  If
all these components dominate \(E\), the image is the generic point of
\(E\), which is not the closed point \(z\); otherwise some \(E_j\) is
contracted to a closed point \(w\neq z\) by hypothesis, and the image is
\(w\).  In both cases the center of \(x\) is not \(z\), so
\(\varphi_z(x)=0\) by \cref{itm:ch22-bump:support}.

\cref{itm:ch22-bump:centers}  The special fibre \(\mathcal Y_s\) has finitely
many components, and each has a well-defined center on \(\mathcal X\).
\end{proof}

\subsection{No countable cofinal system of models}

\begin{example}[The model system is not countably cofinal]
\label{ex:ch22-no-countable-cofinal}
Keep the notation of \Cref{lem:ch22-bump} and assume in addition that
\(k\) is uncountable, for instance \(k=\mathbf C\).  Then \(E(k)\) is
uncountable.  Let \(\mathcal Y_1,\mathcal Y_2,\ldots\) be any sequence of
models of \(X\) over \(R\).  By
\Cref{lem:ch22-bump}\,\cref{itm:ch22-bump:centers} the set

\[
 Z:=\bigcup_{m\geq1}\{z\in E(k):
      z\text{ is the center of a component of }\mathcal Y_{m,s}\}
\]

is countable, so we may choose \(z\in E(k)\setminus Z\).  No
\(\mathcal Y_m\) dominates \(\mathcal X_z\): a birational morphism
\(\mathcal Y_m\to\mathcal X_z\) of models would carry the strict transform
of \(F_z\) to a component of \(\mathcal Y_{m,s}\) with center \(z\) on
\(\mathcal X\).  Hence \(\{\mathcal Y_m\}\) is not cofinal, and the directed
system of models of \(X\) has no countable cofinal subsystem.  Compare
\Cref{ex:no-cofinal-sequence}, of which this is a geometric instance: the
relevant directed set is the set of finite subsets of the uncountable set
\(E(k)\).
\end{example}

\subsection{A compactness argument in which a sequence does not suffice}

\begin{example}[A countable blind spot]
\label{ex:ch22-countable-blind-spot}
Keep the hypotheses of \Cref{ex:ch22-no-countable-cofinal}.

\emph{Two psh functions that countably many skeleta cannot separate.}  Let
\(\mathcal Y_1,\mathcal Y_2,\ldots\) be any prescribed sequence of models
and choose \(z\in E(k)\setminus Z\) as above.  By
\Cref{lem:ch22-bump}\,\cref{itm:ch22-bump:invisible},

\[
 \varphi_z|_{\Delta_{\mathcal Y_m}}=0=0|_{\Delta_{\mathcal Y_m}}
 \quad\text{for every }m,
 \qquad\text{while}\qquad
 \|\varphi_z-0\|_{\infty}=1 .
\tag{22.70}\label{eq:ch22-blind-spot}
\]

Thus the countable family of coordinate maps
\(\varphi\mapsto\varphi|_{\Delta_{\mathcal Y_m}}\) does not separate the
points \(0\) and \(\varphi_z\) of \(\PSH_0(X,\theta)\); the topology it
induces is not Hausdorff, and limits taken in it are not unique.

\emph{Consequence for diagonal extraction.}  The standard proof of
sequential compactness on a countable product would proceed as follows:
given a sequence \((\psi_j)\) in \(\PSH_0(X,\theta)\), use
\Cref{cor:ch22-arzela-ascoli} on \(\Delta_{\mathcal Y_1}\) to extract a
subsequence converging uniformly there, then a further subsequence on
\(\Delta_{\mathcal Y_2}\), and so on, and finally take the diagonal
subsequence.  Every step of this argument is correct, and its conclusion is
correct: the diagonal subsequence converges uniformly on each
\(\Delta_{\mathcal Y_m}\).  What the argument does not do is identify a
limit in \(\PSH_0(X,\theta)\): by \eqref{eq:ch22-blind-spot} the limit is
determined only up to the ambiguity exhibited there, since \(0\) and
\(\varphi_z\) are indistinguishable by the data collected.  The missing
information is exactly the uncountably many coordinates that the sequence
\((\mathcal Y_m)\) omits, and by
\Cref{ex:ch22-no-countable-cofinal} no choice of the sequence removes the
omission.

\emph{A net that converges while staying at distance one.}  Let \(S=E(k)\),
let \(\mathcal A(S)\) be the set of finite subsets of \(S\) directed by
inclusion, and for \(F\in\mathcal A(S)\) choose any \(z(F)\in S\setminus F\)
and put \(u_F:=\varphi_{z(F)}\).  Given a model \(\mathcal Y\), let
\(F_0\subset S\) be the finite set of centers of components of
\(\mathcal Y_s\); for \(F\supset F_0\) we have \(z(F)\notin F_0\), hence
\(u_F|_{\Delta_{\mathcal Y}}=0\).  Therefore

\[
 u_F\longrightarrow0
 \quad\text{in the BFJ topology},
 \qquad\text{while}\qquad
 \|u_F\|_\infty=1\ \text{ for every }F .
\tag{22.71}\label{eq:ch22-net-at-distance-one}
\]

This net has no subsequence in any useful sense: its index set
\(\mathcal A(S)\) has no cofinal sequence by
\Cref{ex:no-cofinal-sequence}, so a subnet indexed by \(\mathbf N\) does not
exist.  It also shows that the BFJ topology is strictly coarser than
uniform convergence on \(X^{\An}\), as announced in
\Cref{rem:ch22-topology-comparison}.
\end{example}

\begin{warning}{What has and has not been proved about sequences}
\Cref{ex:ch22-countable-blind-spot} shows that the compactness
\emph{argument} of \Cref{thm:ch22-proper-normalization} cannot be run with
sequences, and that no countable family of dual complexes determines a
member of \(\PSH(X,\theta)\).  It does not show that
\(\PSH_0(X,\theta)\) fails to be sequentially compact, and we assert
nothing of the kind.  Nor does it forbid sequential statements proved by
other means: \Cref{cor:psh-sequential-regularization} does produce, for each
\(\varphi\in\PSH(X,\theta)\), a decreasing \emph{sequence} of model psh
functions converging to it pointwise, using the Monge--Amp\`ere capacity
estimates of \Cref{chap:bfj-capacity}, which are unavailable at this stage
of the theory and which do not produce a countable cofinal model system
either.  The correct summary is: nets are what the present arguments
deliver; sequences, when available, are theorems with their own proofs.
\end{warning}

\begin{remark}[Where sequences are legitimate in this chapter]
\label{rem:ch22-legitimate-sequences}
Three places, and no others.  On one fixed dual complex, which is compact
metrizable, so that \Cref{cor:ch22-arzela-ascoli} and the Choquet selection
lemma \Cref{lem:choquet-selection} apply; this is the setting of
\Cref{prop:ch23-fixed-skeleton-selection} and of
\Cref{prop:ch22-fixed-skeleton-decreasing} below, and it is the reason
\Cref{rem:choquet-where-countable} lists dual complexes first.  Along the
divisibility net of multiplier ideals, which has the countable cofinal chain
\((m!)_{m\geq1}\); this is \Cref{lem:ch22-divisibility-net}.  And in the
presence of a continuous limit, where Dini's theorem converts a net into a
sequence; this is \Cref{cor:ch22-continuous-decreasing-sequence}.  A
checklist of the general net-versus-sequence traps is
\Cref{rem:net-vs-sequence-checklist}.
\end{remark}

\section{Multiplier ideals on models: the level actually used}
\label{sec:ch22-multiplier}

Decreasing regularization is deduced in
\Cref{sec:ch22-regularization} from one analytic input, the continuity of
psh envelopes (\Cref{thm:bfj-envelope-regularity}), and that input is proved
in \Cref{chap:bfj-envelopes} from asymptotic multiplier ideals on models.
This section states the ideal-theoretic package at the exact level at which
the main text consumes it, records which parts of
\Cref{app:ideals} it corresponds to, and isolates the two formal
consequences that are used here.  We do not reprove the package, and we do
not restate the appendix in a generality it has not been given.

\subsection{The setting on a model}

Let \(\mathcal X\) be a regular projective flat \(R\)-model of \(X\).  A
\emph{vertical ideal} on \(\mathcal X\) is a coherent ideal sheaf
\(\mathfrak a\subset\cO_{\mathcal X}\) whose cosupport is contained in the
special fibre; equivalently \(\mathfrak a\cdot\cO_X=\cO_X\).  Its associated
model function is \(\log|\mathfrak a|\), computed on the normalized blowup as
in \eqref{eq:ch20-ideal-divisor}, and

\[
 \log|\mathfrak a\mathfrak b|=\log|\mathfrak a|+\log|\mathfrak b|,
 \qquad
 \mathfrak a\subset\mathfrak b
 \ \Longrightarrow\
 \log|\mathfrak a|\leq\log|\mathfrak b|\leq0 .
\tag{22.72}\label{eq:ch22-ideal-function-calculus}
\]

A graded sequence of vertical ideals is a family
\(\mathfrak a_\bullet=(\mathfrak a_m)_{m\geq1}\) with
\(\mathfrak a_m\mathfrak a_\ell\subset\mathfrak a_{m+\ell}\), as in
\Cref{def:appC-graded-sequence}; the base ideals
\eqref{eq:ch23-base-ideal} of a model line bundle form the example that
matters.

\Cref{app:ideals} develops asymptotic multiplier ideals for a nonzero graded
sequence on a variety that is \emph{smooth over a field of characteristic
zero}, and generalized test ideals on a regular \(F\)-finite scheme of
characteristic \(p>0\).  A model \(\mathcal X\) is neither: it is regular
and of finite type over the excellent complete discrete valuation ring
\(R\).  The statements used in \Cref{chap:bfj-envelopes} are therefore
imported in the model form proved by
\cite[Appendix~B]{BFJ16}, whose inputs are log resolutions of regular
excellent schemes with residue characteristic zero, relative Nadel
vanishing, and relative Castelnuovo--Mumford regularity over the affine base
\(\Spec R\).  The appendix statements are the corresponding facts over a
field, recorded there with their own proofs and sources; we do not assert
that they apply verbatim on \(\mathcal X\).

\subsection{The package}

\begin{quote}
\textbf{(M)} \emph{Asymptotic multiplier ideals on a regular model, residue
characteristic zero.}  Let \(\mathcal X\) be a regular projective flat
\(R\)-model with \(K\) of residue characteristic zero, and let
\(\mathfrak a_\bullet\) be a nonzero graded sequence of vertical ideals on
\(\mathcal X\).  There are vertical ideals
\(\mathcal J_m:=\mathcal J(\mathfrak a_\bullet^m)\) such that:
\begin{enumerate}
\item[(M1)] \(\mathfrak a_m\subset\mathcal J_m\) for every \(m\geq1\);
\item[(M2)] \(\mathcal J_{m\ell}\subset\mathcal J_m^{\,\ell}\) for all
\(m,\ell\geq1\);
\item[(M3)] there is an ample line bundle \(\mathcal B\) on
\(\mathcal X\), depending only on \(\mathcal X\), such that for every line
bundle \(\mathcal M\) with \(\mathcal M^{\otimes m}\otimes\mathfrak a_m\)
globally generated at sufficiently divisible levels, the sheaf
\(\mathcal B\otimes\mathcal M^{\otimes m}\otimes\mathcal J_m\) is globally
generated for every \(m\);
\item[(M4)] \(v(\mathcal J_m)\geq m\,v(\mathfrak a_\bullet)-A(v)\) for every
divisorial valuation \(v\) with center on \(\mathcal X\), where \(A(v)\) is
the log discrepancy \eqref{eq:appC-log-discrepancy}.
\end{enumerate}
\end{quote}

Properties (M1)--(M3) are \cite[Theorem~B.8 and Remark~B.9]{BFJ16}; they are
exactly the statements quoted before
\Cref{thm:ch23-base-ideal-approximation}, and their variety analogues are
\Cref{prop:appC-asymptotic-multiplier-properties} and
\Cref{thm:appC-uniform-generation-multiplier}.  Property (M4) is the
valuative estimate \eqref{eq:appC-asymptotic-discrepancy}; it is not needed
for the BFJ envelope theorem, which uses the sandwich
\eqref{eq:ch23-uniform-sandwich} instead, but it is the property that the
positive-characteristic substitutes reproduce, and it is recorded here so
that \Cref{sec:ch22-char-zero} can compare the two engines.

\subsection{Two formal consequences}

The first consequence explains why the package is stated for
\emph{asymptotic} ideals: it is what converts subadditivity into
monotonicity, and it is the only monotonicity in the whole regularization
chain that is genuinely sequential.

\begin{lemma}
\label{lem:ch22-divisibility-net}
Assume \textup{(M1)} and \textup{(M2)} and put
\(u_m:=m^{-1}\log|\mathcal J_m|\) and
\(\varphi_m:=m^{-1}\log|\mathfrak a_m|\).  Then:
\begin{enumerate}
\item\label{itm:ch22-div:mono} \(u_{m\ell}\leq u_m\) for all
\(m,\ell\geq1\); that is, \((u_m)\) is a decreasing net for the directed set
\((\mathbf N_{\geq1},\mid)\) ordered by divisibility;
\item\label{itm:ch22-div:above} \(\varphi_m\leq u_m\leq0\) for every \(m\);
\item\label{itm:ch22-div:cofinal} the chain \(m\mapsto m!\) is cofinal in
\((\mathbf N_{\geq1},\mid)\), so the net in
\cref{itm:ch22-div:mono} may be replaced by the decreasing sequence
\((u_{m!})_{m\geq1}\) without loss of information.
\end{enumerate}
\end{lemma}

\begin{proof}
By (M2) and \eqref{eq:ch22-ideal-function-calculus},

\[
 \log|\mathcal J_{m\ell}|\leq\log\bigl|\mathcal J_m^{\,\ell}\bigr|
 =\ell\log|\mathcal J_m| ,
\]

and dividing by \(m\ell\) gives \cref{itm:ch22-div:mono}.  Assertion
\cref{itm:ch22-div:above} is (M1) together with
\eqref{eq:ch22-ideal-function-calculus} and
\(\mathcal J_m\subset\cO_{\mathcal X}\).  For
\cref{itm:ch22-div:cofinal}, every \(m\) divides \(m!\), and
\(m!\mid(m+1)!\), so the chain is increasing and cofinal; a decreasing net
indexed by a directed set with a cofinal chain has the same infimum along
that chain, by \Cref{lem:monotone-net-converges}.
\end{proof}

\begin{remark}[Two different directed sets]
\label{rem:ch22-two-directed-sets}
It is worth pausing on the contrast between
\Cref{lem:ch22-divisibility-net}\,\cref{itm:ch22-div:cofinal} and
\Cref{ex:ch22-no-countable-cofinal}.  Both the divisibility order on
\(\mathbf N\) and the domination order on models are directed sets that are
not totally ordered, and in both cases the theory is naturally indexed by
them.  Only one of them has a countable cofinal chain.  A reader who has
seen the multiplier-ideal argument, where passing from the divisibility net
to the sequence \((m!)\) is harmless, may transfer the habit to the model
system, where it is not.  This is the most common route to a false
sequential statement in this subject.
\end{remark}

The second consequence is the interface with the present chapter.

\begin{proposition}
\label{prop:ch22-multiplier-interface}
The only use of the package \textup{(M)} in
\Cref{chap:bfj-compactness} is through
\Cref{thm:bfj-envelope-regularity}: for every \(f\in C^0(X^{\An})\) the
envelope \(P_\theta(f)\) is a uniform limit on \(X^{\An}\) of
\(\theta\)-psh model functions.  No statement of
\Crefrange{sec:ch22-vertex}{sec:ch22-nets} uses a multiplier ideal, and the deduction
of \Cref{thm:bfj-decreasing-regularization} in
\Cref{sec:ch22-regularization} uses no property of the approximating
functions beyond being \(\theta\)-psh model functions.
\end{proposition}

\begin{proof}
An inspection of the proofs.  The estimates of
\Cref{sec:ch22-vertex,sec:ch22-slope} use relative intersection theory,
\Cref{lem:ch22-negativity}, and \Cref{thm:skeleton-subdivision}; the
topological statements of
\Cref{sec:ch22-topology} use those estimates, Arzel\`a--Ascoli, and
Tychonoff; the examples of \Cref{sec:ch22-nets} use only blowups of closed
points.  In \Cref{sec:ch22-regularization},
\Cref{thm:bfj-envelope-regularity} is quoted twice --- once to produce a
model psh majorant close to \(\varphi\) at a prescribed point, once to
produce a common successor in the directed set of majorants --- and
everything else is order theory, upper semicontinuity, and Dini's theorem.
\end{proof}

\begin{remark}[The role of the fixed twist]
\label{rem:ch22-fixed-twist}
Property (M3) is the one that cannot be weakened to a twist growing with
\(m\).  In the base-ideal argument
\eqref{eq:ch23-uniform-sandwich} the model function of the twisting bundle
enters divided by \(m\); a bundle \(\mathcal B_m\) growing linearly in
\(m\) would contribute a fixed error and destroy the uniformity of the
convergence, which is precisely the point of
\Cref{thm:ch23-base-ideal-approximation}.  The same remark, in the
trivially valued setting and for test ideals, is the checkpoint after
\Cref{thm:appC-uniform-generation-multiplier}.  Uniformity, not
monotonicity, is what the deep input supplies: monotonicity along
divisibility is already given by
\Cref{lem:ch22-divisibility-net}, which used nothing but (M1) and (M2).
\end{remark}

\section{Envelopes and monotone regularization}
\label{sec:ch22-regularization}

The proof of regularization follows the source order

\[
 \begin{aligned}
  &\text{upper-envelope stability}
    \Longrightarrow \text{continuous psh envelopes},\\
  &\text{continuous psh envelopes}
    \Longrightarrow \text{model approximation},\\
  &\text{model approximation}
    \Longrightarrow \text{decreasing model regularization}.
 \end{aligned}
\tag{22.73}\label{eq:ch22-regularization-chain}
\]

The first two arrows are developed in \Cref{chap:bfj-envelopes}.  We state
the final theorem here and prove the order-theoretic last arrow, isolating
the algebraic input as in \Cref{prop:ch22-multiplier-interface}.

\subsection{Continuous majorants of an upper semicontinuous function}

The first arrow of \eqref{eq:ch22-regularization-chain} is applied to
continuous obstacles, but the function to be regularized is only upper
semicontinuous.  The passage from one to the other is the following
elementary lemma, which we prove because it is used three times below and
because its hypotheses --- compact Hausdorff, not compact metrizable --- are
exactly the ones available on \(X^{\An}\).

\begin{lemma}
\label{lem:ch22-continuous-majorant}
Let \(Y\) be a compact Hausdorff space and let
\(f:Y\to[-\infty,+\infty)\) be upper semicontinuous.  Then:
\begin{enumerate}
\item\label{itm:ch22-maj:point} for every \(y\in Y\) and every real
\(c>f(y)\) there is \(g\in C^0(Y)\) with \(g\geq f\) on \(Y\) and
\(g(y)=c\);
\item\label{itm:ch22-maj:inf} the set \(\mathcal C(f)\) of continuous
majorants of \(f\) is nonempty and closed under finite minima, and
\[
 f(y)=\inf\{g(y):g\in\mathcal C(f)\}
 \qquad\text{for every }y\in Y;
\tag{22.74}\label{eq:ch22-continuous-majorant-inf}
\]
in particular \(f\) is the limit of a decreasing net of continuous
functions.
\end{enumerate}
\end{lemma}

\begin{proof}
Upper semicontinuity on a compact space bounds \(f\) above; set
\(M:=\sup_Yf\in[-\infty,+\infty)\) and, given \(y\) and \(c>f(y)\), put
\(M':=\max\{M,c\}\in\mathbf R\).  The set \(K:=\{f\geq c\}\) is closed
because \(f\) is upper semicontinuous, and \(y\notin K\).  If
\(K=\varnothing\), the constant \(g\equiv c\) satisfies
\cref{itm:ch22-maj:point}.  Otherwise, a compact Hausdorff space is normal,
so Urysohn's lemma provides \(h\in C^0(Y,[0,1])\) with \(h(y)=0\) and
\(h\equiv1\) on \(K\).  Put \(g:=c+(M'-c)h\).  On \(K\) we have
\(g=M'\geq M\geq f\); off \(K\) we have \(f<c\leq g\) because \(h\geq0\) and
\(M'\geq c\).  Finally \(g(y)=c\).

For \cref{itm:ch22-maj:inf}, \(\mathcal C(f)\) contains the constant \(M'\)
for any \(M'\geq M\) real, and the minimum of two continuous majorants is a
continuous majorant, so \(\mathcal C(f)\) is nonempty and downward directed.
Every \(g\in\mathcal C(f)\) satisfies \(g(y)\geq f(y)\), and
\cref{itm:ch22-maj:point} produces majorants with value at \(y\) arbitrarily
close to \(f(y)\) from above, or arbitrarily negative when
\(f(y)=-\infty\).  This is \eqref{eq:ch22-continuous-majorant-inf}, and the
net indexed by \(\mathcal C(f)\) with the reverse pointwise order is
decreasing with pointwise limit \(f\), by
\Cref{lem:monotone-net-converges}.
\end{proof}

\begin{remark}[Why not a sequence here either]
On a compact metrizable space the net of
\Cref{lem:ch22-continuous-majorant}\,\cref{itm:ch22-maj:inf} can be replaced
by a decreasing sequence, by choosing majorants adapted to a countable base.
The Berkovich space \(X^{\An}\) is metrizable when \(K\) has a countable
dense subfield, and need not be otherwise; see
\Cref{rem:choquet-where-countable}.  Since nothing below requires a
sequence at this step, we do not impose the hypothesis.
\end{remark}

\subsection{The regularization theorem}

\begin{theorem}\label{thm:bfj-decreasing-regularization}
For every \(\varphi\in\PSH(X,\theta)\), there is a decreasing net

\[
              \varphi_i\in\mathcal D(X)\cap\PSH(X,\theta),
              \qquad
              \varphi_i\searrow\varphi
              \quad\text{pointwise on }X^{\An}.
\tag{22.75}\label{eq:ch22-decreasing-regularization}
\]

If \(\varphi\) is continuous, it is a uniform limit of a sequence of
\(\theta\)-psh model functions.
\end{theorem}

\begin{proof}
Throughout we use the envelope-regularity theorem
\Cref{thm:bfj-envelope-regularity}: for each continuous obstacle \(u\) on
\(X^{\An}\), the envelope \(P_\theta(u)\) is a uniform limit of
\(\theta\)-psh model functions.  Its hypotheses are the standing ones; its
proof, in \Cref{chap:bfj-envelopes}, uses the package \textup{(M)} of
\Cref{sec:ch22-multiplier}.

\paragraph{Step 1: pointwise model majorants.}
Fix \(x\in X^{\An}\) and \(\varepsilon>0\), and first assume
\(\varphi(x)>-\infty\).  By
\Cref{lem:ch22-continuous-majorant}\,\cref{itm:ch22-maj:point} applied on the
compact Hausdorff space \(X^{\An}\), there is \(u\in C^0(X^{\An})\) with
\(u\geq\varphi\) and \(u(x)\leq\varphi(x)+\varepsilon\).  Since \(\varphi\)
is a competitor in \eqref{eq:ch23-envelope}, we have
\(\varphi\leq P_\theta(u)\leq u\).  Choose a \(\theta\)-psh model function
\(\psi_0\) with \(\|\psi_0-P_\theta(u)\|_\infty\leq\varepsilon\) and put
\(\psi:=\psi_0+\varepsilon\).  Then \(\psi\geq P_\theta(u)\geq\varphi\) and

\[
                         \psi(x)\leq\varphi(x)+3\varepsilon.
\tag{22.76}\label{eq:ch22-pointwise-majorant}
\]

If \(\varphi(x)=-\infty\), the same majorant lemma gives, for each
\(A>0\), a continuous \(u_A\geq\varphi\) with \(u_A(x)=-A\); the preceding
construction then produces a \(\theta\)-psh model function
\(\psi_A\geq\varphi\) with \(\psi_A(x)\leq-A+1\).  In both cases

\[
 \varphi(x)=\inf\{\psi(x):
       \psi\in\mathcal D(X)\cap\PSH(X,\theta),\ \psi\geq\varphi\}.
\tag{22.77}\label{eq:ch22-inf-majorants}
\]

\paragraph{Step 2: the index set and its directedness.}
Let

\[
 I:=\{\psi\in\mathcal D(X)\cap\PSH(X,\theta):
       \psi>\varphi\ \text{everywhere on }X^{\An}\},
\tag{22.78}\label{eq:ch22-index-set}
\]

ordered by the reverse pointwise order, so that \(\psi\preceq\psi'\) means
\(\psi\geq\psi'\).  It is nonempty: if \(\psi\) is as in Step 1 then
\(\psi+\delta\in I\) for every \(\delta>0\), because model functions are
finite.

It is directed.  Let \(\psi_1,\psi_2\in I\).  The function
\(\min\{\psi_1,\psi_2\}-\varphi\) is lower semicontinuous with values in
\((0,+\infty]\), being the minimum of two differences of a continuous
function and an upper semicontinuous one; a lower semicontinuous function
attains its infimum on the compact space \(X^{\An}\), so there is
\(\varepsilon>0\) with

\[
                   \min\{\psi_1,\psi_2\}\geq\varphi+3\varepsilon .
\tag{22.79}\label{eq:ch22-strict-gap}
\]

The obstacle \(\min\{\psi_1,\psi_2\}\) is continuous, and
\(\varphi+3\varepsilon\) is \(\theta\)-psh by
\Cref{prop:ch21-elementary-stability}; hence
\(P_\theta(\min\{\psi_1,\psi_2\})\geq\varphi+3\varepsilon\).  Choose a
\(\theta\)-psh model function \(\eta\) within \(\varepsilon\) of that
envelope in the uniform norm.  Then

\[
       \varphi+\varepsilon
       \leq\eta-\varepsilon
       \leq\min\{\psi_1,\psi_2\},
\tag{22.80}\label{eq:ch22-directed-majorant}
\]

so \(\eta-\varepsilon\in I\) dominates both \(\psi_1\) and \(\psi_2\) in the
order of \(I\); that is, it is a common successor.

\paragraph{Step 3: the net.}
Consider the net indexed by \(I\) whose value at \(\psi\) is \(\psi\)
itself.  It is decreasing by construction.  Its pointwise infimum is
\(\varphi\): the inequality \(\geq\varphi\) is the definition of \(I\), and
for the reverse inequality, given \(c>\varphi(x)\), Step 1 produces
\(\psi\geq\varphi\) with \(\psi(x)<c\), and then
\(\psi+\delta\in I\) with \((\psi+\delta)(x)<c\) for small \(\delta>0\).
By \Cref{lem:monotone-net-converges} the net converges pointwise to its
infimum, which proves \eqref{eq:ch22-decreasing-regularization}.  Note that
\Cref{thm:ch21-decreasing-stability} is consistent with this: the limit is
\(\varphi\), which is psh by hypothesis and in particular not identically
\(-\infty\).

\paragraph{Step 4: the continuous case.}
If \(\varphi\) is continuous, the decreasing net of continuous functions
\((\psi)_{\psi\in I}\) has continuous pointwise limit \(\varphi\) on the
compact space \(X^{\An}\), so Dini's theorem for nets,
\Cref{prop:dini-nets}, gives uniform convergence.  For each \(m\) there is
therefore \(\psi\in I\) with \(\|\psi-\varphi\|_\infty<1/m\), and
\Cref{cor:ch22-continuous-decreasing-sequence} below organizes these choices
into a decreasing sequence.  In particular \(\varphi\) is a uniform limit of
a sequence of \(\theta\)-psh model functions.
\end{proof}

\begin{corollary}
\label{cor:ch22-continuous-decreasing-sequence}
Let \(\varphi\in\PSH(X,\theta)\) be continuous.  Then there is a
\emph{decreasing} sequence \((\varphi_m)_{m\geq1}\) of \(\theta\)-psh model
functions with
\(\|\varphi_m-\varphi\|_{\infty}<1/m\) for every \(m\); in particular
\(\varphi_m\searrow\varphi\) uniformly on \(X^{\An}\).
\end{corollary}

\begin{proof}
Let \(I\) be as in \eqref{eq:ch22-index-set}.  By Step 4 of the preceding
proof, for every \(m\) the set
\(I_m:=\{\psi\in I:\|\psi-\varphi\|_\infty<1/m\}\) is nonempty.  Choose
\(\varphi_1\in I_1\).  Given \(\varphi_{m}\), pick \(\eta\in I_{m+1}\) and,
using directedness of \(I\), a common successor
\(\varphi_{m+1}\in I\) of \(\varphi_m\) and \(\eta\); then
\(\varphi_{m+1}\leq\min\{\varphi_m,\eta\}\) and
\(\varphi_{m+1}\geq\varphi\), so
\(\|\varphi_{m+1}-\varphi\|_\infty\leq\|\eta-\varphi\|_\infty<1/(m+1)\).
The resulting sequence is decreasing and converges uniformly.
\end{proof}

\begin{remark}[An improvement on the naive selection]
\label{rem:ch22-decreasing-selection-improvement}
A weaker and more common formulation of the continuous case selects, for
each \(m\), some member of the regularizing net within \(1/m\) of
\(\varphi\), and then observes that the selected members converge uniformly
without claiming that they decrease.  \Cref{cor:ch22-continuous-decreasing-sequence}
shows that directedness of \(I\) allows the selection to be made decreasing,
at no cost.  This matters where monotonicity is used later --- for instance
in the monotone extension of the Monge--Amp\`ere operator
(\Cref{thm:bounded-ma-extension}) --- and it is the form in which
\Cref{thm:ch21-continuous-zhang} uses the present theorem.  What is
\emph{not} improved is the singular case: there the index set \(I\) has no
countable cofinal subset in general, and no argument in this chapter
produces a sequence.
\end{remark}

\subsection{Decreasing selection on a fixed skeleton}

On one dual complex the situation is different, because a dual complex is
compact metrizable.  The following principle is the decreasing counterpart
of the fixed-skeleton selection \Cref{prop:ch23-fixed-skeleton-selection},
and it is proved by the same dense-set, common-modulus and Dini
architecture.

\begin{proposition}
\label{prop:ch22-fixed-skeleton-decreasing}
Let \(\mathcal X\) be an SNC model determining \(\theta\), write
\(\Delta=\Delta_{\mathcal X}\), and let
\(\mathcal G\subset\PSH(X,\theta)\) be nonempty and downward directed for
the pointwise order on \(X^{\An}\).  Assume that

\[
 u:=\inf_{\psi\in\mathcal G}\psi|_\Delta
\]

is finite at one point of \(\Delta\).  Then \(u\) is finite and continuous
on \(\Delta\), and there is a decreasing sequence
\((\psi_j)_{j\geq1}\) in \(\mathcal G\) with
\(\psi_j|_\Delta\to u\) uniformly on \(\Delta\).
\end{proposition}

\begin{proof}
\emph{Common modulus.}  Fix \(\psi_0\in\mathcal G\) and discard from
\(\mathcal G\) every member not dominated by \(\psi_0\); by directedness the
remaining family is still downward directed with the same infimum.  All
members now satisfy \(\psi\leq\psi_0\), so their suprema are bounded above,
and by hypothesis \(u(y_0)>-\infty\) for some \(y_0\in\Delta\), whence
\(\psi(y_0)\geq u(y_0)\) for all \(\psi\).  Applying
\Cref{thm:bfj-skeletal-equicontinuity} to
\(\psi-\sup\psi\) and \Cref{cor:ch22-explicit-modulus}, we obtain one
constant \(C\) such that every \(\psi|_\Delta\) is \(C\)-Lipschitz for
\(d_\Delta\); combined with the two-sided bound at \(y_0\) and the bounded
diameter of \(\Delta\), the restrictions are uniformly bounded on
\(\Delta\).  A pointwise infimum of a family of \(C\)-Lipschitz functions
with a common bound is \(C\)-Lipschitz and finite; hence \(u\) is continuous
on \(\Delta\).

\emph{Dense set and finite minima.}  The complex \(\Delta\) is compact
metrizable by \Cref{lem:ch22-path-metric}; choose a countable dense set
\((x_k)_{k\geq1}\).  For all \(k,m\geq1\) select
\(\eta_{k,m}\in\mathcal G\) with

\[
                     \eta_{k,m}(x_k)<u(x_k)+m^{-1},
\tag{22.81}\label{eq:ch22-dense-selection}
\]

which is possible by definition of the infimum.  Enumerate the countable
family \(\{\eta_{k,m}\}\) as \((\eta_j)_{j\geq1}\) and construct
\(\psi_j\in\mathcal G\) recursively: let \(\psi_1\) be a common successor of
\(\psi_0\) and \(\eta_1\), and let \(\psi_{j+1}\) be a common successor of
\(\psi_j\) and \(\eta_{j+1}\).  Directedness makes this possible, and by
construction \((\psi_j)\) is decreasing with
\(\psi_j\leq\eta_j\).

\emph{Conclusion.}  The decreasing sequence \((\psi_j|_\Delta)\) has a
pointwise limit \(w\geq u\) on \(\Delta\), and \(w\leq\eta_j\) forces
\(w(x_k)\leq u(x_k)\) for every \(k\) by
\eqref{eq:ch22-dense-selection}; hence \(w=u\) on a dense set.  Both
functions are \(C\)-Lipschitz, so \(w=u\) on \(\Delta\).  Since \(u\) is
continuous and the convergence is monotone on a compact space, Dini's
theorem \Cref{prop:dini-nets} makes it uniform.
\end{proof}

\begin{warning}{The selection is skeleton by skeleton}
\Cref{prop:ch22-fixed-skeleton-decreasing} produces one sequence for one
prescribed dual complex.  It asserts nothing about all complexes at once,
and the corresponding global statement is not available: it would require a
countable cofinal system of models, which does not exist by
\Cref{ex:ch22-no-countable-cofinal}.  The same restriction is imposed on the
increasing version \Cref{prop:ch23-fixed-skeleton-selection}, and for the
same reason.  Whenever a later argument needs a sequence on a skeleton --- a
measure carried by one complex, an integral against such a measure --- it
must name the complex first.
\end{warning}

\begin{remark}[Hypotheses of \Cref{thm:bfj-decreasing-regularization}, in full view]
\label{rem:ch22-regularization-hypotheses}
The theorem has been proved under the standing hypotheses, and every one of
them is used.  The field is complete discretely valued of residue
characteristic zero: discreteness enters through the model theory of
\Cref{chap:bfj-model-forms} and residue characteristic zero through
\Cref{prop:snc-cofinal} and the package \textup{(M)}.  The variety is smooth
and projective: smoothness supports resolution and the SNC models, and
projectivity supports the ample class and Serre's theorems.  The class
\(\{\theta\}\) is ample: this is what makes
\(\PSH(X,\theta)\) nonempty (\Cref{prop:ch20-positive-representative}) and
what the envelope theorem assumes.  No algebraizability of \(X\) over a
function field is assumed here; that hypothesis appears in the original
solution of the Monge--Amp\`ere equation and is removed in
\Cref{thm:bggjkm-remove-algebraizability}, for a different reason ---
orthogonality, not regularization.  Finally, the conclusion is a net; the
sequential strengthening for continuous \(\varphi\) is
\Cref{cor:ch22-continuous-decreasing-sequence}, and the sequential
strengthening in general is \Cref{cor:psh-sequential-regularization}, proved
by capacity in \Cref{chap:bfj-capacity}.
\end{remark}

\section{The exact role of residue characteristic zero}
\label{sec:ch22-char-zero}

Residue characteristic zero has been invoked three times in this chapter,
always by reference and never inside a computation.  Because it is easy to
attribute more to it than is true, and equally easy to attribute less, this
section is an inventory: what uses it, what does not, what is known without
it, and what must not be said.

\subsection{Inventory}

\begin{center}
\small
\begin{tabular}{p{0.30\textwidth}p{0.24\textwidth}p{0.36\textwidth}}
\toprule
Statement & Uses residue char.\ \(0\)? & Through what\\
\midrule
\Cref{lem:ch22-incidence-signs} & no & intersection theory, formal
functions\\
\Cref{prop:ch22-vertex-estimate} & no & the above and connectedness\\
\Cref{lem:ch22-nef-approximation}, \Cref{lem:ch22-negativity} & no &
Kleiman's criterion, relative generation\\
\Cref{lem:ch22-star-shrinking} & no & the toroidal part of
\Cref{thm:skeleton-subdivision}\\
\Cref{prop:ch22-slope-estimate} & no & the four items above\\
\Cref{thm:bfj-skeletal-equicontinuity} & no & the two estimates\\
\Cref{def:ch22-bfj-topology} & \emph{yes} &
\Cref{prop:snc-cofinal}: the coordinates must be cofinal\\
\Cref{thm:ch22-proper-normalization}, \Cref{thm:bfj-psh-compactness}
 & \emph{yes}, indirectly & through the definition of the topology and
   \Cref{thm:bfj-skeletal-characterization}\\
Package \textup{(M)}, \Cref{sec:ch22-multiplier} & \emph{yes} &
log resolutions, subadditivity, relative Nadel vanishing\\
\Cref{thm:bfj-decreasing-regularization} & \emph{yes} & only through
\Cref{thm:bfj-envelope-regularity}\\
\Cref{prop:ch22-fixed-skeleton-decreasing},
\Cref{cor:ch22-continuous-decreasing-sequence} & no & order theory and
Dini, given the input\\
\bottomrule
\end{tabular}
\end{center}

The first six rows are \Cref{prop:ch22-characteristic-free-estimates}.  The
remaining ones split into two genuinely different uses, discussed next.

\subsection{First use: cofinality of SNC models}

The definition of a BFJ psh function quantifies over SNC models
(\Cref{def:ch21-bfj-psh}), the reconstruction formula
\eqref{eq:ch21-skeletal-reconstruction} takes an infimum over them, and the
topology \eqref{eq:ch22-product-embedding} takes a product over them.  All
three are informative only because SNC models are cofinal among all models,
which is \Cref{prop:snc-cofinal}, itself an application of resolution and
principalization in the form \Cref{thm:appB-resolution}.  If cofinality
fails, the union of the dual complexes of those models that happen to be SNC
need not be dense in \(X^{\An}\), the inverse-limit description
\Cref{thm:skeleton-inverse-limit} may fail, and the coordinates of
\eqref{eq:ch22-product-embedding} may lose information --- an uncountable
version of the blind spot of
\Cref{ex:ch22-countable-blind-spot}.  This is spelled out in
\Cref{rem:cofinality-scope}, which is the place in this book where the
failure is analyzed.

Two further consequences in the immediate neighbourhood should be recorded.
The exactness of \eqref{eq:ch20-ddc-exact} was proved on regular models and
then passed to the limit using their cofinality, so it too carries the
hypothesis.  And the substitute usually proposed --- de Jong's semistable
alteration theorem, \Cref{thm:appB-de-jong} --- produces a generically
finite cover rather than a birational modification, hence does not restore
cofinality of models of the given \(X\); it changes the variety.

\subsection{Second use: the multiplier-ideal package}

The package \textup{(M)} of \Cref{sec:ch22-multiplier} uses residue
characteristic zero in each of its three parts, and for different reasons:

\begin{enumerate}
\item the definition of \(\mathcal J(\mathfrak a_\bullet^m)\) requires a log
resolution of the pair, \Cref{def:appB-log-resolution}, which is available
for excellent regular schemes with residue characteristic zero by
\Cref{thm:appB-resolution}, and is not known in general otherwise;
\item subadditivity (M2) is the theorem of Demailly, Ein and Lazarsfeld
\Cref{thm:appC-multiplier-subadditivity}, whose proof restricts to the
diagonal of a product and uses smoothness there
\cite[Theorem~2.3]{DEL00};
\item uniform global generation (M3) is proved from relative Nadel vanishing
and Castelnuovo--Mumford regularity, \Cref{thm:appC-uniform-generation-multiplier};
Kodaira-type vanishing theorems are false in positive characteristic, so
this proof does not transpose.
\end{enumerate}

By contrast, (M4) is a valuative estimate, and it does have a
positive-characteristic analogue: the test-ideal discrepancy inequality
\eqref{eq:appC-test-discrepancy}.  This asymmetry is the reason the
positive-characteristic theory of \Cref{chap:bfj-continuity} is possible at
all, and also the reason it is not a formal transcription: the objects
change, and only their inequalities survive.

\subsection{What is known without residue characteristic zero}

\begin{enumerate}
\item \emph{The estimates.}  \Cref{prop:ch22-characteristic-free-estimates}:
vertex and slope bounds, hence facewise equicontinuity, hold on any SNC
model over any complete discretely valued field.  What is lost is not the
estimate but the supply of complexes.
\item \emph{Curves.}  For an ample line bundle on a smooth projective
geometrically connected curve over an arbitrary complete nontrivially
valued field, the semipositive metric envelope is continuous
\cite[Theorem~1.1]{GJKM19}; see
\Cref{thm:curves-arbitrary-field-continuity}.  This is a rational-model
envelope statement.  Transporting
\Cref{thm:bfj-decreasing-regularization} to that setting additionally
requires its singular-psh class, skeletal hypotheses, and envelope
identification; it is not a consequence of notation alone.
\item \emph{Equal positive characteristic.}  Over a complete discretely
valued field of characteristic \(p>0\), for a smooth projective variety of
geometric origin from a \(d\)-dimensional family, and assuming resolution
and embedded resolution over the base field in dimension \(d+n\),
Gubler--Jell--K\"unnemann--Martin prove continuity of envelopes with test
ideals in place of multiplier ideals and deduce monotone regularization;
this is \Cref{thm:gjkm-positive-characteristic-cy} and \cite{GJKM19}.  The
hypotheses are unconditional for surfaces of the stated geometric origin.
\item \emph{Mixed characteristic.}  Conditionally on resolution for
projective models and embedded resolution for regular projective models,
the same architecture runs with perturbation-friendly test ideals; this is
\Cref{thm:fgk-mixed-characteristic-cy}.
\item \emph{The trivially valued theory.}  \Cref{part:bj} proves the
envelope property in arbitrary characteristic by a different route, using
test ideals on \(X\times\mathbf A^1\) in characteristic \(p\)
(\Cref{chap:bj-smooth-characteristic}) and finite descent; it is not
obtained from the present chapter by specialization.
\end{enumerate}

\begin{warning}{Four statements that would be false or unproved}
\begin{enumerate}
\item ``SNC models are cofinal'' --- unproved in positive or mixed residue
characteristic in dimension \(\geq3\); see \Cref{rem:cofinality-scope}.
\item ``Multiplier ideals may be defined on any regular model'' --- the
definition \eqref{eq:appC-multiplier-resolution} needs a log resolution.
\item ``\Cref{thm:bfj-decreasing-regularization} holds over any complete
discretely valued field'' --- what holds in that generality is
\Cref{prop:ch22-characteristic-free-estimates}, not the theorem.
\item ``Test ideals are multiplier ideals in characteristic \(p\)'' --- the
comparison \Cref{thm:appC-hara-yoshida} holds only after spreading out and
excluding finitely many primes, and
\Cref{ex:ch22-fpt-cusp} shows that the two invariants genuinely differ.
\end{enumerate}
\end{warning}

\subsection{The boundary, computed on an example}

The failure in positive characteristic is not a counterexample to
compactness or to regularization; no such counterexample is known.  It is
the failure of the \emph{inputs}, and the following classical computation
shows that the failure is real and quantitative rather than a defect of
proof technique: the characteristic-\(p\) substitute for a multiplier ideal
computes a different number.

\begin{example}[The cusp: a threshold that depends on \(p\)]
\label{ex:ch22-fpt-cusp}
Let \(f=x^2+y^3\) in two variables over an algebraically closed field.  In
characteristic zero the log canonical threshold at the origin is

\[
 \lct(f)=\frac12+\frac13=\frac56 ,
\tag{22.82}\label{eq:ch22-lct-cusp}
\]

the value at which the multiplier ideal
\(\mathcal J((f)^c)\) first becomes nontrivial; this is
\Cref{ex:appC-snc} applied to the log resolution of the cusp.  In
characteristic \(p>0\) the corresponding invariant for test ideals is the
\(F\)-pure threshold
\(\operatorname{fpt}(f)=\lim_{e\to\infty}p^{-e}\nu(p^e)\), where
\(\nu(p^e)\) is the largest \(r\) with
\(f^r\notin(x^{p^e},y^{p^e})\).  Expanding
\(f^r=\sum_k\binom rkx^{2(r-k)}y^{3k}\) and applying Lucas's congruence to
the binomial coefficients gives

\[
 \operatorname{fpt}(x^2+y^3)=
 \begin{cases}
  \tfrac12,&p=2,\\[1mm]
  \tfrac23,&p=3,\\[1mm]
  \tfrac56,&p\equiv1\ (\mathrm{mod}\ 6),\\[1mm]
  \tfrac56-\dfrac1{6p},&p\equiv5\ (\mathrm{mod}\ 6).
 \end{cases}
\tag{22.83}\label{eq:ch22-fpt-cusp}
\]

Three sample verifications, each a finite computation.  For \(p=3\) and
\(e=2\) the admissible pairs \((r,k)\) satisfy \(k\leq2\) and
\(r-k\leq4\); the only candidate with \(r=6\) is \(k=2\), and
\(\binom62=15\equiv0\), while \(r=5,k=1\) gives \(\binom51=5\equiv2\), so
\(\nu(9)=5\) and \(5/9=\tfrac23-\tfrac19\).  For \(p=5\) and \(e=2\) one
finds \(\nu(25)=19\) and \(19/25=\tfrac45-\tfrac1{25}\), consistent with
\(\operatorname{fpt}=\tfrac45=\tfrac56-\tfrac1{30}\).  For \(p=7\) one finds
\(\nu(7^e)=5(7^e-1)/6\), so \(\operatorname{fpt}=\tfrac56\) exactly, as
predicted by \(6\mid p-1\).

Two conclusions.  First, for the infinitely many primes
\(p\equiv5\ (\mathrm{mod}\ 6)\) the characteristic-\(p\) invariant is
strictly smaller than \eqref{eq:ch22-lct-cusp}, so no characteristic-free
identity of ideals can hold; this is the quantitative content of the
firewall \Cref{check:appC-characteristic-firewall}, and it is compatible
with \Cref{thm:appC-hara-yoshida} because the discrepancy
\(1/(6p)\) tends to zero, so equality is recovered after excluding finitely
many primes for each fixed datum.  Second, what survives uniformly in \(p\)
is the \emph{inequality} \eqref{eq:appC-test-discrepancy}, and it is
precisely the inequalities, not the identities, that the envelope argument
of \cite{GJKM19} consumes.  The general theory of these thresholds is in
\cite{HY03} and \cite{ST12}; the computation above is the standard cusp
example.
\end{example}

\begin{remark}[How to state a theorem in this chapter]
\label{rem:ch22-how-to-state}
The discipline that follows from this section is simple.  An estimate proved
by intersection theory on a fixed model may be stated for a complete
discretely valued field of arbitrary residue characteristic.  A statement
that quantifies over all SNC models, or that uses the reconstruction of a
psh function from its skeletal restrictions, carries residue characteristic
zero.  A statement whose proof passes through envelopes carries it twice
over.  In the rest of \Cref{part:bfj}, the standing hypotheses of
\Cref{chap:bfj-model-forms} are in force, so the distinction is invisible;
it becomes visible again in
\Cref{chap:bfj-continuity}, whose hypothesis atlas should be read together
with the present inventory.
\end{remark}

\section{Worked examples}
\label{sec:ch22-worked-examples}

\begin{example}[A decreasing regularization computed on a curve skeleton]
\label{ex:ch22-regularization-curve}
We compute the net of \Cref{thm:bfj-decreasing-regularization} explicitly in
a case where it can be written down.

\emph{The data.}  Let \(k\) be algebraically closed of characteristic zero,
\(R=k[[t]]\), \(K=k((t))\), \(X=\mathbf P^1_K\) with coordinate \(w\), and
\(\mathcal X=\mathbf P^1_R\), so that \(\mathcal X_s=E\cong\mathbf P^1_k\) is
reduced and irreducible and \(\Delta_{\mathcal X}\) is the single Gauss point
\(x_E\).  Let \(L=\cO_X(1)\), let \(\theta\) be the curvature of the model
metric induced by \(\cO_{\mathcal X}(1)\), so \((\theta\cdot E)=\deg L=1\),
and put

\[
 \varphi(x):=-\max\{0,v_x(w)\}
 =\min\{0,\log|w(x)|\},
 \qquad x\in X^{\An}.
\tag{22.84}\label{eq:ch22-green-function}
\]

Thus \(\varphi=0\) outside the open disc \(\{|w|<1\}\), \(\varphi(x)=-v_x(w)\)
inside it, and \(\varphi=-\infty\) exactly at the \(K\)-point \(w=0\).

\emph{The models.}  Let \(\mathcal X_m\) be obtained from \(\mathcal X\) by
blowing up \(m\) times, starting at the reduction of the section
\(\{w=0\}\) and continuing at the intersection of the last exceptional
divisor with the strict transform of that section.  All exceptional
components \(F_1,\ldots,F_m\) have multiplicity one, and

\[
 v_{x_{F_j}}(w)=j,\qquad j=1,\ldots,m .
\]

The dual complex \(\Delta_{\mathcal X_m}\) is the chain
\(x_E,x_{F_1},\ldots,x_{F_m}\); in the coordinate \(s=v(w)\) it is the
segment \([0,m]\) with vertices at the integers, each edge of lattice length
one, as in \Cref{ex:dual-complex-chain-cycle}.

\emph{The regularizing sequence.}  Put

\[
 \varphi_m:=\max\{\varphi,-m\}
 =\varphi_{D_m},
 \qquad
 D_m:=-\sum_{j=1}^{m}jF_j .
\tag{22.85}\label{eq:ch22-truncated-green}
\]

The second equality holds because a model function determined on
\(\mathcal X_m\) is constant on the fibres of \(p_{\mathcal X_m}\) and
because the coefficients of \(D_m\) give, by
\eqref{eq:ch20-sign}, the values \(0,-1,\ldots,-m\) at the vertices; on
\(X^{\An}\) both sides equal \(-\min\{v_x(w),m\}\) for
\(v_x(w)\geq0\) and \(0\) otherwise.

\emph{The nef test, with numbers.}  Write \(\rho_m:\mathcal X_m\to\mathcal X\)
for the composed blowup.  The reference degrees on
\(\mathcal X_m\) are \((\rho_m^*\theta\cdot E)=1\) and
\((\rho_m^*\theta\cdot F_j)=0\), and the incidence graph is a chain, so
\eqref{eq:ch20-curve-nef-test} reads

\[
 \begin{aligned}
 \text{at }E:&\quad 1+(-1-0)=0\geq0,\\
 \text{at }F_j,\ 1\leq j\leq m-1:&\quad
   0+\bigl(-(j-1)+j\bigr)+\bigl(-(j+1)+j\bigr)=0\geq0,\\
 \text{at }F_m:&\quad 0+\bigl(-(m-1)+m\bigr)=1\geq0 .
 \end{aligned}
\tag{22.86}\label{eq:ch22-green-nef-test}
\]

Hence each \(\varphi_m\) is a \(\theta\)-psh model function, and
\(\varphi_m\searrow\varphi\) pointwise.  This is an explicit instance of
\eqref{eq:ch22-decreasing-regularization}, computed without any appeal to
envelopes or multiplier ideals.

\emph{The slope is exactly critical.}  The first line of
\eqref{eq:ch22-green-nef-test} vanishes: the degree \((\theta\cdot E)=1\) is
exactly consumed by the outgoing slope \(-1\).  Consequently
\(c\varphi\) is \emph{not} \(\theta\)-psh for \(c>1\): the same test gives
\(1-c<0\).  The maximal admissible depth rate is a numerical invariant of
the polarization, not a defect of the construction; for
\(L=\cO_X(d)\) the admissible rate is \(d\).

\emph{Why a sequence sufficed here.}  The polar set of \(\varphi\) is one
\(K\)-point, and the models \(\mathcal X_m\) form a countable chain adapted
to it.  Nothing of the kind is available in general: by
\Cref{ex:ch22-no-countable-cofinal} the model system has no countable
cofinal subsystem, and by \Cref{prop:ch22-multiplier-interface} the general
proof produces only a net.  A general sequential statement does exist, but
its proof is the capacity argument of
\Cref{cor:psh-sequential-regularization}, not the argument of this chapter.
\end{example}

\begin{example}[Why one skeleton is not enough]
\label{ex:ch22-one-skeleton-insufficient}
Blow up a closed point \(z\) of one component of a regular model, as in
\Cref{lem:ch22-bump}.  The old dual complex embeds in the refined one by
\Cref{prop:skeleton-nesting}, but the exceptional component contributes a
new vertex.  A vertical divisor supported on that exceptional component,
small enough to be absorbed by the fixed relatively ample reference class,
has a \(\theta\)-psh model function that vanishes at every old vertex and is
nonzero at the new one: this is exactly
\(\varphi_z\) of \eqref{eq:ch22-bump-function}, with

\[
 \varphi_z|_{\Delta_{\mathcal X}}=0,
 \qquad
 \varphi_z(x_{F_z})=-1 .
\]

Thus convergence on a single skeleton cannot define the BFJ topology; every
fixed skeleton is only one finite-dimensional coordinate of
\eqref{eq:ch22-product-embedding}, and by
\Cref{ex:ch22-countable-blind-spot} countably many coordinates are still not
enough.
\end{example}

\begin{example}[Dini selection]
\label{ex:ch22-dini-selection}
Suppose a decreasing regularizing net \(\varphi_i\) has continuous limit
\(\varphi\).  The open sets

\[
                     U_i^{(m)}=\{\varphi_i<\varphi+1/m\}
\]

increase with \(i\) and cover the compact Berkovich space \(X^{\An}\).
Directedness gives one index \(i(m)\) with
\(U_{i(m)}^{(m)}=X^{\An}\).  Hence

\[
             0\leq\varphi_{i(m)}-\varphi<1/m
             \quad\text{uniformly}.
\]

This is the proof of \Cref{prop:dini-nets} written out in the case at hand,
and it produces the sequence in the last assertion of
\Cref{thm:bfj-decreasing-regularization} without assuming that the model
system is countable.  Combining it with directedness, as in
\Cref{cor:ch22-continuous-decreasing-sequence}, makes the selected sequence
decreasing as well.
\end{example}

\begin{example}[Compactness of a normalized family, tested on the segment]
\label{ex:ch22-compactness-segment}
Return to \Cref{ex:ch22-modulus-segment}, with \(d_0=1\), \(d_1=2\).  By
\eqref{eq:ch22-curve-slope-form} and
\Cref{thm:bfj-skeletal-equicontinuity}, the normalized restrictions to the
segment all lie in

\[
 \mathcal K=\bigl\{u:[0,1]\to[-2,0]\ \text{convex},\
       u'(0^+)\geq-1,\ u'(1^-)\leq2,\ \max u=0\bigr\},
\]

a uniformly bounded, uniformly Lipschitz, closed subset of
\(C^0([0,1])\), hence compact by Arzel\`a--Ascoli, in accordance with
\Cref{cor:ch22-arzela-ascoli}.  Its deepest member is
\eqref{eq:ch22-flat-extremal}, with minimum \(-2\): if the maximum is
attained at \(t=1\), convexity and \(u'(1^-)\leq2\) give
\(u(t)\geq2(t-1)\), and if it is attained at \(t=0\), convexity and
\(u'(0^+)\geq-1\) give \(u(t)\geq-t\).  Compactness on
this one complex is elementary and metric.  What
\Cref{thm:bfj-psh-compactness} adds is the compatibility of all such
compactness statements across the uncountable system of models, and that is
where Tychonoff and the net formulation enter.
\end{example}

\section{Where the deep inputs occur}
\label{sec:ch22-ledger}

The following ledger prevents later arguments from attributing too much to
compactness alone.  It refines the table of
\Cref{sec:ch22-signature} by naming, for each step, the nonformal
ingredient and the statement in which it is proved or imported.

\begin{center}
\small
\begin{tabular}{p{0.26\textwidth}p{0.64\textwidth}}
\toprule
Step & Nonformal input \\
\midrule
Vertex bounds & Relative nef intersection and connectedness of the special
fibre; see \eqref{eq:ch22-vertex-intersection} and
\Cref{lem:ch22-incidence-signs}.\\
Nef approximation & Kleiman's criterion in the relative form,
\Cref{lem:ch22-nef-approximation}; relative global generation.\\
Negativity lemma & Vertical \(\mathbf Q\)-factoriality and relative
generation, \Cref{lem:ch22-negativity}; no vanishing theorem.\\
Slope bounds & Projective toroidal subdivisions
(\Cref{lem:ch22-star-shrinking}), their realization by vertical blowups
(\Cref{thm:skeleton-subdivision}), and the negativity lemma; see
\eqref{eq:ch22-slope-intersection}.\\
Simplex lemma & Convex geometry only, \Cref{lem:ch22-simplex}; the constant
is \eqref{eq:ch22-gauge}.\\
Compactness & The preceding uniform estimates; Arzel\`a--Ascoli and
Tychonoff are then formal topology, and the index set is uncountable
(\Cref{ex:ch22-no-countable-cofinal}).\\
The topology is informative & Cofinality of SNC models,
\Cref{prop:snc-cofinal}, hence resolution of singularities.\\
Upper-envelope stability & Skeletal equicontinuity and convexity, in
\Cref{thm:ch23-upper-envelope}; no multiplier ideal is used.\\
Envelope regularity & Log resolutions, asymptotic multiplier ideals,
subadditivity, and uniform global generation on regular models: the package
\textup{(M)} of \Cref{sec:ch22-multiplier}.\\
Decreasing net & The elementary order construction
\eqref{eq:ch22-inf-majorants}--\eqref{eq:ch22-directed-majorant}, once
envelope regularity is known, together with
\Cref{lem:ch22-continuous-majorant}.\\
Decreasing sequence, continuous case & Dini for nets and directedness,
\Cref{cor:ch22-continuous-decreasing-sequence}; no countability of the model
system.\\
Decreasing sequence, general case & Not proved here: it is the capacity
argument \Cref{cor:psh-sequential-regularization}.\\
\bottomrule
\end{tabular}
\end{center}

Residue characteristic zero is used to obtain the cofinal regular and SNC
models and the multiplier-ideal inputs in the original proof, as itemized in
\Cref{sec:ch22-char-zero}.  Arzel\`a--Ascoli, Tychonoff, Dini, and the final
order argument have no characteristic restriction of their own.

\section{Chapter dictionary}
\label{sec:ch22-dictionary}

\[
\begin{array}{ccl}
\text{relative nefness on a model}
 &\longleftrightarrow&
 \text{linear inequalities on vertex values},\\[1mm]
\text{toroidal star shrinking}
 &\longleftrightarrow&
 \text{derivatives become vertex values},\\[1mm]
\text{uniform intersection bounds}
 &\longleftrightarrow&
 \text{equicontinuity on }\Delta_{\mathcal X},\\[1mm]
\text{product over models}
 &\longleftrightarrow&
 \text{BFJ topology; nets, not sequences},\\[1mm]
\text{continuity of envelopes}
 &\longleftrightarrow&
 \text{decreasing model regularization}.
\end{array}
\]

In more detail, the chapter establishes the following correspondences and
records where each is used next.

\begin{center}
\small
\begin{tabular}{p{0.27\textwidth}|p{0.34\textwidth}|p{0.31\textwidth}}
\textbf{Object}&\textbf{Established property}&\textbf{First use}\\ \hline
incidence numbers \(w_{ij}\)
 &Laplacian signs \eqref{eq:ch22-off-diagonal-positive},
  \eqref{eq:ch22-row-sum-zero}
 &\Cref{prop:ch22-vertex-estimate}\\
vertex estimate
 &explicit constant \eqref{eq:ch22-vertex-lower-bound}
 &\Cref{prop:ch22-slope-estimate}, \Cref{chap:bfj-energy}\\
star shrinking
 &\Cref{lem:ch22-star-shrinking}, realized by a
  \(\mathbf Q\)-factorial toroidal blowup
 &\Cref{prop:ch22-slope-estimate}\\
negativity lemma
 &\eqref{eq:ch22-negativity}
 &\Cref{prop:ch22-slope-estimate}\\
simplex lemma
 &\eqref{eq:ch22-simplex-lipschitz} with gauge
  \eqref{eq:ch22-gauge}
 &\Cref{prop:ch22-slope-estimate},
  \Cref{ex:ch22-modulus-triangle}\\
skeletal estimate
 &\Cref{thm:bfj-skeletal-equicontinuity}
 &\Cref{thm:ch21-decreasing-stability},
  \Cref{thm:ch23-upper-envelope},
  \Cref{chap:bfj-finite-energy-measures}\\
modulus of continuity
 &\eqref{eq:ch22-modulus}
 &\Cref{prop:ch23-fixed-skeleton-selection},
  \Cref{prop:ch22-fixed-skeleton-decreasing}\\
BFJ topology
 &\Cref{def:ch22-bfj-topology}; three equivalent descriptions
 &\Cref{thm:bfj-psh-compactness}, \Cref{chap:bfj-energy}\\
properness of \(\sup\)
 &\Cref{thm:ch22-proper-normalization}
 &\Cref{thm:bfj-psh-compactness}\\
compactness modulo constants
 &\Cref{thm:bfj-psh-compactness}, in net form
 &\Cref{chap:bfj-finite-energy-measures},
  \Cref{chap:bfj-continuity}\\
continuous majorants
 &\Cref{lem:ch22-continuous-majorant}
 &\Cref{thm:bfj-decreasing-regularization}\\
decreasing regularization
 &\Cref{thm:bfj-decreasing-regularization}
 &\Cref{thm:ch21-continuous-zhang},
  \Cref{thm:bounded-ma-extension},
  \Cref{chap:bfj-capacity}\\
decreasing sequence, continuous case
 &\Cref{cor:ch22-continuous-decreasing-sequence}
 &\Cref{thm:ch21-continuous-zhang}\\
characteristic ledger
 &\Cref{prop:ch22-characteristic-free-estimates},
  \Cref{sec:ch22-char-zero}
 &\Cref{chap:bfj-continuity}\\
\end{tabular}
\end{center}

\section*{Notes and further reading}

The facewise Lipschitz estimate is \cite[Theorem~6.1]{BFJ16}; its passage to
general psh functions, the properness of the normalization, compactness, and
the equivalence of the three descriptions of the topology are in
\cite[Section~7]{BFJ16}.  Decreasing regularization is
\cite[Theorem~8.7]{BFJ16}.  Its multiplier-ideal input is presented in
\Cref{chap:bfj-envelopes}, with the general algebra of asymptotic multiplier
ideals reviewed in \Cref{app:ideals} and the model-level package quoted in
\Cref{sec:ch22-multiplier} from \cite[Appendix~B]{BFJ16}.  The negativity
lemma in the form used here is \cite[Lemma~1.6]{BFJ16}; the realization of
projective subdivisions by normalized blowups is
\cite[Theorem~3.15]{BFJ16}, proved in this book as
\Cref{thm:skeleton-subdivision}.

Three points of presentation depart from the source and should be flagged.
First, the exponent in the slope--intersection inequality
\eqref{eq:ch22-slope-intersection} is \(n-\#J\); the printed display
\cite[(6.3)]{BFJ16} has \(n-\#J-1\), which leaves a one-cycle rather than a
number and is negative when \(n=\#J=1\).  The correction is dimensional, it
changes no sign or normalization, and it is recorded with its verification
in the editorial ledger and in \Cref{app:provenance}.  Second, the simplex
lemma is proved here with an explicit constant, the gauge
\eqref{eq:ch22-gauge} of the convex hull of the edge vectors, rather than
quoted; the gauge is computed in
\Cref{ex:ch22-gauge-computed} and shown to be attained in
\Cref{ex:ch22-modulus-triangle}.  Third, the continuous majorant lemma
\Cref{lem:ch22-continuous-majorant}, which is what allows an upper
semicontinuous \(\varphi\) to be attacked by continuous obstacles, is proved
here rather than assumed; on a compact Hausdorff space it is a two-line
consequence of Urysohn's lemma, and it is the only place where normality of
\(X^{\An}\) is used.

The present proof deliberately keeps nets visible.  \Cref{app:nets} records
the compactness, Dini, and Choquet facts in the form used, and
\Cref{chap:topological-tools} proves them.  The blind-spot example
\Cref{ex:ch22-countable-blind-spot} is elementary and is included because
the temptation to extract a diagonal subsequence over countably many
skeleta is strong and the failure is invisible in the notation; the
corresponding general warning is
\Cref{rem:net-vs-sequence-checklist}.  The fixed-skeleton selection
principles, in their increasing form
\Cref{prop:ch23-fixed-skeleton-selection} and their decreasing form
\Cref{prop:ch22-fixed-skeleton-decreasing}, are the sanctioned way to
reintroduce sequences; both are restricted to one compact dual complex by
design.

For the positive-characteristic and mixed-characteristic landscape, and for
the exact hypotheses under which the analogues of the results of this
chapter are known, see \Cref{chap:bfj-continuity} and its hypothesis atlas;
the primary sources are \cite{GJKM19} for equal characteristic \(p\) with
test ideals, \cite{BGGJKM20} for the removal of algebraizability in residue
characteristic zero, and \cite{GM19} for the comparison of algebraic,
formal, and piecewise linear semipositivity over more general fields.  The
threshold computation \Cref{ex:ch22-fpt-cusp} is the standard cusp example
of the test-ideal literature; see \cite{HY03} and \cite{ST12}.
